\documentclass[12pt,a4paper]{article}
\usepackage[utf8]{inputenc}
\usepackage[T1]{fontenc}
\usepackage[spanish]{babel}
\usepackage[margin=2.5cm]{geometry}
\usepackage{graphicx}
\usepackage{longtable}
\usepackage{array}
\usepackage{hyperref}

\graphicspath{{../03_Modelado/Diagramas_UML/}}

\title{Especificación de Requisitos de Software\\Secciones 2 y 3: Bloque individual\\[2mm]\large Versión 2.0 (PE5): cierre del ERS con los componentes de inteligencia artificial}
\author{Calle Delgado Kamila Anabella}
\date{23 de agosto de 2026}

\begin{document}
\maketitle
\setcounter{section}{1}

% Seccion 2: Descripcion general
\section{Descripción general}
\label{sec:descripcion-general}

% Subseccion 2.1: Perspectiva del producto y limite del sistema
\subsection{Perspectiva del producto y límite del sistema}
\label{sec:perspectiva-producto}

El Sistema de Gestión Agrícola con Inteligencia Artificial para Cultivos de Verde y Cacao (SGA) es una aplicación orientada a centralizar el registro y la consulta de la información generada por las actividades agrícolas de la finca, sustituyendo el bloc de datos y la bitácora en papel que describe el administrador (EV-01 [03:34]). El sistema administra la información correspondiente a parcelas, cultivos, trabajadores, labores, cosechas e inventario de insumos, e incorpora un módulo de Inteligencia Artificial que analiza los registros históricos para generar alertas y recomendaciones relacionadas con el manejo de los cultivos.

La finca donde se levantó la evidencia de campo mantiene un policultivo de palma, verde y cacao (EV-01 [03:18]). El administrador señala que la labor que más le interesaría automatizar es la cosecha de palma africana (EV-01 [07:28]); no obstante, el alcance de este proyecto se limita a los cultivos de verde y cacao. Esta delimitación es una decisión de alcance del proyecto: la gestión de palma queda fuera del sistema en la presente versión.

El límite del sistema separa lo que el SGA administra internamente de las entidades externas con las que intercambia información. El módulo de Inteligencia Artificial se representa dentro de ese límite, ya que la Sección 5 no evidencia que dependa de ningún servicio de terceros.

% Subseccion 2.1.1: Diagrama de contexto
\subsubsection{Diagrama de contexto}
\label{sec:diagrama-contexto-texto}

La Figura~\ref{fig:diagrama-contexto} muestra el diagrama de contexto del SGA. Hacia el sistema, el Administrador registra parcelas, trabajadores, cultivos, actividades, inventario y cosechas; el sistema le devuelve reportes, alertas, recomendaciones de IA y estadísticas. El Jornalero recibe del sistema las tareas, la parcela asignada y las actividades, y registra hacia el sistema las labores realizadas, la cosecha y las incidencias. El Centro de acopio se incorpora como entidad externa adicional respecto del diagrama de la Entrega 1B: el peso y el comprobante que entrega respaldan el registro de cada cosecha (EV-01 [04:32]). Las otras dos entidades candidatas evaluadas, proveedor de insumos (EV-01 [05:50]) y contador externo (EV-01 [06:23], [06:37]), no muestran intercambio de datos con el sistema y no se incorporan.

\begin{figure}[tbp]
\centering
\includegraphics[width=0.85\textwidth,height=0.4\textheight,keepaspectratio]{diagrama_contexto.png}
\caption{Diagrama de contexto del SGA}
\label{fig:diagrama-contexto}
\end{figure}

% Subseccion 2.2: Mapa de partes interesadas con matriz poder/interes
\subsection{Mapa de partes interesadas con matriz poder/interés}
\label{sec:mapa-stakeholders}

La identificación de partes interesadas del SGA parte de la Tabla A.1 del Anexo A del ERS de la Entrega 1B, que reconoce cuatro actores: el administrador de la finca, el jornalero, el equipo de desarrollo y el docente de la asignatura.

% Subseccion 2.2.1: Identificacion de partes interesadas
\subsubsection{Identificación de partes interesadas}
\label{sec:identificacion-stakeholders}

\begin{center}
\begin{longtable}{|p{1.1cm}|p{2.7cm}|p{2.8cm}|p{1.3cm}|p{1.3cm}|p{2.1cm}|p{1.1cm}|}
\hline
\textbf{ID} & \textbf{Actor} & \textbf{Rol} & \textbf{Poder} & \textbf{Interés} & \textbf{Cuadrante} & \textbf{ID-EV} \\
\hline
\endfirsthead
\hline
\textbf{ID} & \textbf{Actor} & \textbf{Rol} & \textbf{Poder} & \textbf{Interés} & \textbf{Cuadrante} & \textbf{ID-EV} \\
\hline
\endhead
\hline
\endfoot
\hline
\endlastfoot
PI-01 & Administrador de la finca & Cliente principal & Alto & Alto & Gestionar de cerca & EV-01 \\
\hline
PI-02 & Jornalero & Usuario final & Medio-bajo & Bajo & Monitorear & EV-02 \\
\hline
PI-03 & Equipo de desarrollo & Desarrolladores & Alto & Alto & Gestionar de cerca & N/D \\
\hline
PI-04 & Docente de la asignatura & Supervisor académico & Medio & Medio & Cruce, sin inclinación & N/D \\
\hline
\end{longtable}
\label{tab:identificacion-stakeholders}
\end{center}

% Subseccion 2.2.2: Matriz poder/interes
\subsubsection{Matriz poder/interés}
\label{sec:matriz-poder-interes-texto}

La Figura~\ref{fig:matriz-poder-interes} ubica a los cuatro actores en la matriz poder/interés. El Docente de la asignatura se representa con un marcador distinto (rombo) exactamente en el cruce de los dos ejes, ya que la Tabla A.1 declara su poder e interés como medios, sin evidencia en la Sección 5 que lo incline hacia un cuadrante específico.

% Subseccion 2.2.3: Estrategia de gestion por cuadrante
\subsubsection{Estrategia de gestión por cuadrante}
\label{sec:estrategia-cuadrante}

Para el Administrador de la finca y el Equipo de desarrollo, ubicados en el cuadrante de gestionar de cerca, la estrategia es de comunicación directa y frecuente, con participación en la priorización de requisitos: el administrador es la única parte interesada con entrevista extensa (EV-01, 640,88 s) y el equipo de desarrollo tiene la responsabilidad operativa directa sobre el sistema, declarada en la Tabla A.1.

Para el Jornalero, ubicado hacia el cuadrante de monitorear, la estrategia es de comunicación periódica y de baja carga, limitada a la información mínima necesaria para ejecutar sus tareas: su propia declaración (EV-02 [02:43]) describe una disposición pragmática, no proactiva, hacia el uso de herramientas digitales.

Para el Docente de la asignatura, ubicado en el cruce de los dos ejes, la estrategia es de mantenerlo informado mediante los entregables periódicos del proyecto: su interés declarado en la Tabla A.1 corresponde a una supervisión académica periódica, no a una participación operativa diaria.

\begin{figure}[tbp]
\centering
\includegraphics[width=0.6\textwidth,height=0.4\textheight,keepaspectratio]{mapa_stakeholders_poder_interes.png}
\caption{Matriz poder/interés de las partes interesadas del SGA}
\label{fig:matriz-poder-interes}
\end{figure}

% Subseccion 2.3: Modelado organizacional i estrella
\subsection{Modelado organizacional en el lenguaje distribuido intencional (i*)}
\label{sec:modelado-istar}

El modelado organizacional del SGA utiliza la notación i* de Eric Yu para representar los actores del dominio, sus dependencias mutuas y el razonamiento interno que sostiene cada una de ellas.

% Subseccion 2.3.1: Actores y dependencias del dominio
\subsubsection{Actores y dependencias del dominio}
\label{sec:actores-dependencias}

Los actores del modelo i* coinciden exactamente con los identificados en la Sección~\ref{sec:mapa-stakeholders}: Administrador de la finca, Jornalero, Equipo de desarrollo y Docente de la asignatura. La Sección 5 del material fuente no registra ninguna mención al Equipo de desarrollo ni al Docente, por lo que ambos se representan como actores sin dependencias evidenciadas hacia o desde el dominio operativo de la finca.

Entre el Administrador y el Jornalero se identifican cuatro dependencias, todas ancladas a la evidencia de campo:

\begin{itemize}
\item El Jornalero depende del Administrador para la asignación de la tarea diaria (EV-02 [01:01]).
\item El Jornalero depende del Administrador para la entrega de insumos agrícolas (EV-02 [02:25]).
\item El Administrador depende del Jornalero para el reporte de incidencias fitosanitarias (EV-02 [01:34]).
\item El Administrador depende del Jornalero para el registro de la producción cosechada (EV-01 [05:18]).
\end{itemize}

% Subseccion 2.3.2: Diagrama de Dependencia Estrategica
\subsubsection{Diagrama de Dependencia Estratégica (SD)}
\label{sec:istar-sd-texto}

La Figura~\ref{fig:istar-sd} representa las cuatro dependencias identificadas entre el Administrador y el Jornalero. El Equipo de desarrollo y el Docente de la asignatura se muestran como actores sin enlaces de dependencia, conforme a lo señalado en la Sección~\ref{sec:actores-dependencias}.

\begin{figure}[tbp]
\centering
\includegraphics[width=0.9\textwidth,height=0.4\textheight,keepaspectratio]{istar_sd.png}
\caption{Diagrama de Dependencia Estratégica (SD) del SGA}
\label{fig:istar-sd}
\end{figure}

% Subseccion 2.3.3: Diagrama de Razon Estrategica
\subsubsection{Diagrama de Razón Estratégica (SR)}
\label{sec:istar-sr-texto}

La Figura~\ref{fig:istar-sr} abre el razonamiento interno de cada actor. El Administrador persigue el objetivo duro de gestionar eficientemente la producción de la finca, que se descompone en las tareas de planificar actividades semanalmente (EV-01 [06:51]), asignar tareas a los trabajadores (EV-01 [06:58]) y consultar reportes de producción (EV-01 [08:10]). Declara además dos objetivos blandos: confiar en las recomendaciones de la Inteligencia Artificial (EV-01 [09:02], [09:10]), sostenido por la tarea de verificar manualmente la alerta antes de actuar, y reducir el tiempo dedicado a labores manuales de control (EV-01 [01:36]), para el cual la Sección 5 no evidencia una tarea que contribuya directamente.

El Jornalero persigue el objetivo duro de cumplir las tareas asignadas, descompuesto en ejecutar la labor correspondiente (EV-02 [00:54]) y reportar novedades al administrador (EV-02 [01:34]). Declara el objetivo blando de minimizar el uso de dispositivos nuevos (EV-02 [02:33], [02:43]), sin una tarea evidenciada que contribuya a él de forma específica.

El Equipo de desarrollo persigue el objetivo duro de entregar el SGA conforme a los requisitos definidos, mediante la tarea de analizar, diseñar e implementar el sistema. El Docente de la asignatura persigue el objetivo duro de verificar el cumplimiento de la metodología y los objetivos de la asignatura, mediante la tarea de evaluar los entregables del proyecto.

\begin{figure}[tbp]
\centering
\includegraphics[width=0.7\textwidth,height=0.55\textheight,keepaspectratio]{istar_sr.png}
\caption{Diagrama de Razón Estratégica (SR) del SGA}
\label{fig:istar-sr}
\end{figure}

% Seccion 3: Requisitos especificos completos
\section{Requisitos específicos completos}
\label{sec:requisitos-especificos}

\subsection{Requisitos funcionales}
\label{sec:requisitos-funcionales}

% Subseccion 3.1.1: Plantilla de especificacion
\subsubsection{Plantilla de especificación}
\label{sec:plantilla-rf}

Los requisitos funcionales del Sistema de Gestión Agrícola con Inteligencia Artificial para Cultivos de Verde y Cacao (SGA) se documentan mediante la plantilla de ocho atributos que se describe en la Tabla~\ref{tab:plantilla-rf}. Esta plantilla amplía la utilizada en la Entrega 2 (1B), que incluía Identificador, Nombre, Descripción, Actor(es), Prioridad, Precondiciones, Flujo principal y Postcondiciones, incorporando un origen verificable por evidencia de campo, la especificación de entradas y salidas, y un criterio de verificación cuantificable.

\begin{center}
\begin{longtable}{|p{3.5cm}|p{10.7cm}|}
\hline
\textbf{Atributo} & \textbf{Contenido} \\
\hline
\endfirsthead
\hline
\textbf{Atributo} & \textbf{Contenido} \\
\hline
\endhead
\hline
\endfoot
\hline
\endlastfoot
Identificador & Código único con el formato RF-XX. \\
\hline
Nombre & Frase nominal breve que identifica la capacidad especificada. \\
\hline
Descripción & Enunciado en la forma «El sistema deberá\ldots», limitado a una única capacidad por ficha. \\
\hline
Actor / origen & Actor que origina o utiliza el requisito, junto con el identificador de evidencia EV-XX y la marca de tiempo que lo respalda, o la restricción normativa aplicable cuando el requisito no se origina en una cita directa. \\
\hline
Entradas & Datos que el sistema recibe para ejecutar la capacidad, con su tipo y su obligatoriedad. \\
\hline
Salidas & Datos o efectos que el sistema produce como resultado de la capacidad. \\
\hline
Prioridad (MoSCoW) & Clasificación en Debe tener, Debería tener, Podría tener o No tendrá. \\
\hline
Criterio de verificación & Condición cuantificable, con métrica, umbral y método de comprobación, que permite determinar si el requisito se cumple. \\
\hline
\end{longtable}
\label{tab:plantilla-rf}
\end{center}

% Subseccion 3.1.2: Catalogo de requisitos funcionales
\subsubsection{Catálogo de requisitos funcionales}
\label{sec:catalogo-rf}

El catálogo siguiente incorpora los dieciséis requisitos funcionales identificados en la Entrega 2 (1B), migrados a la plantilla de ocho atributos sin alterar su contenido sustantivo, y su ampliación hasta un total de treinta y un requisitos funcionales derivados de la evidencia de campo descrita en la Sección 5. Los identificadores RF-01 a RF-16 corresponden a los requisitos migrados; los identificadores RF-17 a RF-27 corresponden a los requisitos incorporados en la Entrega 3 (2A); los identificadores RF-28 a RF-31 se incorporan en la práctica PE5 y especifican los componentes de inteligencia artificial del sistema (Sección~\ref{sec:componentes-ia}). En los requisitos migrados se conservan, como campos adicionales al cierre de la ficha, las Precondiciones, el Flujo principal y las Postcondiciones ya redactados en la Entrega 2 (1B).

Los requisitos RF-01 y RF-14 no derivan de una cita textual de las entrevistas, ya que ninguna pregunta del guion abordó de forma directa el mecanismo de autenticación. Su origen se declara como la restricción R-02 sobre acceso diferenciado por rol, necesaria para sostener la distinción de responsabilidades entre Administrador y Trabajador agrícola evidenciada a lo largo de ambas entrevistas.

\begin{minipage}{\linewidth}
\bigskip
\noindent\textbf{RF-01. Inicio de sesión}
\label{tab:RF-01}

\begin{center}
\begin{tabular}{|p{3.5cm}|p{10.7cm}|}
\hline
Identificador & RF-01 \\
\hline
Nombre & Inicio de sesión \\
\hline
Descripción & El sistema deberá permitir que el administrador y los trabajadores agrícolas inicien sesión mediante un usuario y una contraseña para acceder a las funcionalidades correspondientes a su rol. \\
\hline
Actor / origen & Administrador, Trabajador agrícola: R-02 (el acceso a la información dependerá del rol asignado a cada usuario). \\
\hline
Entradas & Usuario (texto, obligatorio); contraseña (texto, obligatorio). \\
\hline
Salidas & Sesión iniciada con acceso a las funciones del rol correspondiente, o mensaje de error. \\
\hline
Prioridad (MoSCoW) & Debe tener. \\
\hline
Criterio de verificación & El sistema autentica al usuario y concede acceso según su rol en $\leq$ 2,0 s, verificado en el percentil 95 de 50 intentos de inicio de sesión. \\
\hline
Precondiciones & El usuario debe estar registrado en el sistema. \\
\hline
Flujo principal & 1. El usuario ingresa sus credenciales. 2. El sistema valida la información. 3. Si las credenciales son correctas, permite el acceso al sistema según el rol del usuario. \\
\hline
Postcondiciones & El usuario accede correctamente al sistema. \\
\hline
\end{tabular}
\end{center}
\end{minipage}

\begin{minipage}{\linewidth}
\bigskip
\noindent\textbf{RF-02. Gestión de parcelas}
\label{tab:RF-02}

\begin{center}
\begin{tabular}{|p{3.5cm}|p{10.7cm}|}
\hline
Identificador & RF-02 \\
\hline
Nombre & Gestión de parcelas \\
\hline
Descripción & El sistema deberá permitir registrar, consultar, actualizar y eliminar la información correspondiente a las parcelas o lotes que conforman la finca. \\
\hline
Actor / origen & Administrador: EV-01 [04:02], [04:14]. \\
\hline
Entradas & Nombre o código de parcela (texto, obligatorio); ubicación (texto, obligatorio); área (numérico, obligatorio); estado (selección, obligatorio). \\
\hline
Salidas & Parcela registrada, actualizada o eliminada, disponible para su asociación con cultivos (RF-03). \\
\hline
Prioridad (MoSCoW) & Debe tener. \\
\hline
Criterio de verificación & La consulta de una parcela devuelve sus campos obligatorios en $\leq$ 2,0 s con 500 parcelas cargadas, verificado en el percentil 95 de 50 ejecuciones. \\
\hline
Precondiciones & El administrador debe haber iniciado sesión. \\
\hline
Flujo principal & 1. El administrador selecciona el módulo de parcelas. 2. Registra o modifica la información requerida. 3. El sistema valida y almacena los datos. \\
\hline
Postcondiciones & La información de la parcela queda registrada o actualizada correctamente. \\
\hline
\end{tabular}
\end{center}
\end{minipage}

\begin{minipage}{\linewidth}
\bigskip
\noindent\textbf{RF-03. Gestión de cultivos}
\label{tab:RF-03}

\begin{center}
\begin{tabular}{|p{3.5cm}|p{10.7cm}|}
\hline
Identificador & RF-03 \\
\hline
Nombre & Gestión de cultivos \\
\hline
Descripción & El sistema deberá permitir registrar y administrar los cultivos de verde y cacao asociados a cada parcela de la finca. \\
\hline
Actor / origen & Administrador: EV-01 [03:05], [03:18]. \\
\hline
Entradas & Parcela asociada (selección, obligatorio); tipo de cultivo (selección: verde o cacao, obligatorio); variedad (texto, opcional); fecha de cultivo (fecha, obligatorio). \\
\hline
Salidas & Cultivo registrado y asociado a la parcela seleccionada. \\
\hline
Prioridad (MoSCoW) & Debe tener. \\
\hline
Criterio de verificación & El registro de un cultivo queda disponible para consulta asociado a su parcela en $\leq$ 2,0 s, verificado sobre 30 registros consecutivos. \\
\hline
Precondiciones & Debe existir al menos una parcela registrada. \\
\hline
Flujo principal & 1. El administrador selecciona una parcela. 2. Registra la información del cultivo. 3. El sistema almacena los datos asociados a la parcela seleccionada. \\
\hline
Postcondiciones & El cultivo queda asociado correctamente a la parcela correspondiente. \\
\hline
\end{tabular}
\end{center}
\end{minipage}

\begin{minipage}{\linewidth}
\bigskip
\noindent\textbf{RF-04. Registro de actividades agrícolas}
\label{tab:RF-04}

\begin{center}
\begin{tabular}{|p{3.5cm}|p{10.7cm}|}
\hline
Identificador & RF-04 \\
\hline
Nombre & Registro de actividades agrícolas \\
\hline
Descripción & El sistema deberá permitir registrar las actividades realizadas por los trabajadores agrícolas, indicando la fecha, la actividad ejecutada y la parcela donde fue desarrollada. \\
\hline
Actor / origen & Trabajador agrícola: EV-02 [00:54]; EV-01 [07:11]. \\
\hline
Entradas & Fecha (fecha, obligatorio); tipo de actividad (selección, obligatorio); parcela (selección, obligatorio). \\
\hline
Salidas & Actividad registrada, disponible para consulta del administrador (RF-10) y para el historial de trabajador (RF-24). \\
\hline
Prioridad (MoSCoW) & Debe tener. \\
\hline
Criterio de verificación & El 100~\% de las actividades registradas por un trabajador quedan disponibles para consulta del administrador en un máximo de 1 minuto, verificado sobre una muestra de 30 registros. \\
\hline
Precondiciones & El trabajador debe haber iniciado sesión. \\
\hline
Flujo principal & 1. El trabajador accede al módulo de actividades. 2. Registra la información solicitada. 3. El sistema valida y almacena el registro. \\
\hline
Postcondiciones & La actividad queda registrada correctamente en el sistema. \\
\hline
\end{tabular}
\end{center}
\end{minipage}

\begin{minipage}{\linewidth}
\bigskip
\noindent\textbf{RF-05. Registro de cosechas}
\label{tab:RF-05}

\begin{center}
\begin{tabular}{|p{3.5cm}|p{10.7cm}|}
\hline
Identificador & RF-05 \\
\hline
Nombre & Registro de cosechas \\
\hline
Descripción & El sistema deberá permitir registrar la producción obtenida en cada cosecha, especificando el cultivo, la parcela, la cantidad recolectada, la unidad de medida y la fecha del registro. \\
\hline
Actor / origen & Trabajador agrícola: EV-01 [04:32], [04:47]; EV-02 [02:08]. \\
\hline
Entradas & Cultivo (selección, obligatorio); parcela (selección, obligatorio); cantidad (numérico, obligatorio); unidad de medida (selección, obligatorio); fecha (fecha, obligatorio). \\
\hline
Salidas & Producción registrada e incorporada al historial de producción. \\
\hline
Prioridad (MoSCoW) & Debe tener. \\
\hline
Criterio de verificación & El registro de una cosecha se incorpora al historial de producción en $\leq$ 2,0 s, verificado en el percentil 95 de 30 ejecuciones. \\
\hline
Precondiciones & Debe existir un cultivo previamente registrado en el sistema. \\
\hline
Flujo principal & 1. El trabajador selecciona el cultivo correspondiente. 2. Ingresa la información de la cosecha. 3. El sistema valida y almacena el registro. \\
\hline
Postcondiciones & La producción queda registrada y disponible para consultas y reportes. \\
\hline
\end{tabular}
\end{center}
\end{minipage}

\begin{minipage}{\linewidth}
\bigskip
\noindent\textbf{RF-06. Gestión de inventario de insumos}
\label{tab:RF-06}

\begin{center}
\begin{tabular}{|p{3.5cm}|p{10.7cm}|}
\hline
Identificador & RF-06 \\
\hline
Nombre & Gestión de inventario de insumos \\
\hline
Descripción & El sistema deberá permitir registrar, consultar, actualizar y eliminar la información correspondiente a los insumos agrícolas y herramientas utilizadas en la finca, incluyendo fertilizantes, fungicidas, herbicidas y herramientas de trabajo. \\
\hline
Actor / origen & Administrador: EV-01 [05:36], [05:44], [05:58], [06:14]. \\
\hline
Entradas & Nombre del insumo (texto, obligatorio); tipo (selección: fertilizante, fungicida, herbicida o herramienta, obligatorio); cantidad disponible (numérico, obligatorio); unidad de medida (selección, obligatorio). \\
\hline
Salidas & Inventario actualizado, disponible para consulta y para el descuento por entregas (RF-22). \\
\hline
Prioridad (MoSCoW) & Debe tener. \\
\hline
Criterio de verificación & El inventario refleja una actualización de cantidad disponible en $\leq$ 2,0 s tras el registro, verificado sobre 30 actualizaciones consecutivas. \\
\hline
Precondiciones & El administrador debe haber iniciado sesión. \\
\hline
Flujo principal & 1. El administrador accede al módulo de inventario. 2. Registra o actualiza la información del insumo. 3. El sistema valida y almacena los cambios realizados. \\
\hline
Postcondiciones & El inventario queda actualizado con la información registrada. \\
\hline
\end{tabular}
\end{center}
\end{minipage}

\begin{minipage}{\linewidth}
\bigskip
\noindent\textbf{RF-07. Consulta de producción por período}
\label{tab:RF-07}

\begin{center}
\begin{tabular}{|p{3.5cm}|p{10.7cm}|}
\hline
Identificador & RF-07 \\
\hline
Nombre & Consulta de producción por período \\
\hline
Descripción & El sistema deberá permitir consultar la producción agrícola registrada durante un período específico, facilitando el seguimiento del rendimiento de los cultivos. \\
\hline
Actor / origen & Administrador: EV-01 [04:47], [04:54]. \\
\hline
Entradas & Fecha de inicio y fecha de fin del período (fecha, obligatorio). \\
\hline
Salidas & Listado de producción registrada dentro del período seleccionado. \\
\hline
Prioridad (MoSCoW) & Debe tener. \\
\hline
Criterio de verificación & La consulta de producción por período con 500 registros cargados responde en $\leq$ 2,0 s, verificado en el percentil 95 de 50 ejecuciones, y presenta únicamente los registros del período seleccionado. \\
\hline
Precondiciones & Deben existir registros de producción almacenados en el sistema. \\
\hline
Flujo principal & 1. El administrador selecciona el período de consulta. 2. El sistema procesa la información registrada. 3. Se presenta el resultado de la consulta. \\
\hline
Postcondiciones & La información solicitada se muestra correctamente al usuario. \\
\hline
\end{tabular}
\end{center}
\end{minipage}

\begin{minipage}{\linewidth}
\bigskip
\noindent\textbf{RF-08. Planificación de actividades agrícolas}
\label{tab:RF-08}

\begin{center}
\begin{tabular}{|p{3.5cm}|p{10.7cm}|}
\hline
Identificador & RF-08 \\
\hline
Nombre & Planificación de actividades agrícolas \\
\hline
Descripción & El sistema deberá permitir planificar las actividades agrícolas indicando la tarea, la fecha programada y la parcela donde será ejecutada. \\
\hline
Actor / origen & Administrador: EV-01 [06:51], [06:58]. \\
\hline
Entradas & Tarea (texto, obligatorio); fecha programada (fecha, obligatorio); parcela (selección, obligatorio). \\
\hline
Salidas & Actividad planificada, disponible para su posterior asignación (RF-09). \\
\hline
Prioridad (MoSCoW) & Debe tener. \\
\hline
Criterio de verificación & Una actividad planificada queda disponible para su asignación en $\leq$ 2,0 s desde su registro, verificado sobre 30 registros consecutivos. \\
\hline
Precondiciones & Deben existir parcelas registradas en el sistema. \\
\hline
Flujo principal & 1. El administrador accede al módulo de planificación. 2. Registra la actividad programada. 3. El sistema almacena la planificación. \\
\hline
Postcondiciones & La actividad queda programada para su posterior ejecución. \\
\hline
\end{tabular}
\end{center}
\end{minipage}

\begin{minipage}{\linewidth}
\bigskip
\noindent\textbf{RF-09. Asignación de tareas}
\label{tab:RF-09}

\begin{center}
\begin{tabular}{|p{3.5cm}|p{10.7cm}|}
\hline
Identificador & RF-09 \\
\hline
Nombre & Asignación de tareas \\
\hline
Descripción & El sistema deberá permitir asignar actividades específicas a los trabajadores agrícolas, indicando la tarea, la parcela y la fecha de ejecución. \\
\hline
Actor / origen & Administrador: EV-01 [06:58]; EV-02 [01:01], [01:18]. \\
\hline
Entradas & Trabajador (selección, obligatorio); actividad (selección, obligatorio); parcela (selección, obligatorio); fecha de ejecución (fecha, obligatorio). \\
\hline
Salidas & Tarea asignada al trabajador seleccionado. \\
\hline
Prioridad (MoSCoW) & Debe tener. \\
\hline
Criterio de verificación & El 100~\% de las tareas asignadas quedan disponibles para consulta del trabajador correspondiente en $\leq$ 2,0 s, verificado sobre una muestra de 30 asignaciones. \\
\hline
Precondiciones & Deben existir trabajadores y actividades registrados en el sistema. \\
\hline
Flujo principal & 1. El administrador selecciona al trabajador. 2. Asigna la actividad correspondiente. 3. El sistema registra la asignación realizada. \\
\hline
Postcondiciones & La tarea queda asignada al trabajador seleccionado. \\
\hline
\end{tabular}
\end{center}
\end{minipage}

\begin{minipage}{\linewidth}
\bigskip
\noindent\textbf{RF-10. Seguimiento del cumplimiento de tareas}
\label{tab:RF-10}

\begin{center}
\begin{tabular}{|p{3.5cm}|p{10.7cm}|}
\hline
Identificador & RF-10 \\
\hline
Nombre & Seguimiento del cumplimiento de tareas \\
\hline
Descripción & El sistema deberá permitir consultar el estado de las actividades asignadas a los trabajadores, identificando aquellas pendientes, en proceso o finalizadas. \\
\hline
Actor / origen & Administrador: EV-01 [07:11], [07:18]. \\
\hline
Entradas & Ninguna directa; opera sobre tareas previamente asignadas (RF-09). \\
\hline
Salidas & Estado de cada tarea asignada (pendiente, en proceso o finalizada). \\
\hline
Prioridad (MoSCoW) & Debe tener. \\
\hline
Criterio de verificación & El estado de una tarea asignada se actualiza y refleja correctamente en $\leq$ 2,0 s desde el cambio registrado, verificado sobre 30 tareas de prueba. \\
\hline
Precondiciones & Deben existir tareas previamente asignadas. \\
\hline
Flujo principal & 1. El administrador accede al módulo de seguimiento. 2. El sistema consulta el estado de las tareas. 3. Se presenta la información correspondiente. \\
\hline
Postcondiciones & El administrador conoce el avance de las actividades programadas. \\
\hline
\end{tabular}
\end{center}
\end{minipage}

\begin{minipage}{\linewidth}
\bigskip
\noindent\textbf{RF-11. Generación de reportes de producción}
\label{tab:RF-11}

\begin{center}
\begin{tabular}{|p{3.5cm}|p{10.7cm}|}
\hline
Identificador & RF-11 \\
\hline
Nombre & Generación de reportes de producción \\
\hline
Descripción & El sistema deberá permitir generar reportes de producción por cultivo, parcela o período, facilitando el análisis de la información registrada. \\
\hline
Actor / origen & Administrador: EV-01 [08:10]. \\
\hline
Entradas & Filtros de consulta: cultivo, parcela o período (selección, al menos uno obligatorio). \\
\hline
Salidas & Reporte de producción según los filtros seleccionados, disponible para consulta o descarga. \\
\hline
Prioridad (MoSCoW) & Debe tener. \\
\hline
Criterio de verificación & El sistema genera el reporte de producción solicitado en $\leq$ 3,0 s con 500 registros cargados, verificado en el percentil 95 de 30 ejecuciones. \\
\hline
Precondiciones & Deben existir registros de producción almacenados en el sistema. \\
\hline
Flujo principal & 1. El administrador accede al módulo de reportes. 2. Selecciona los filtros de consulta. 3. El sistema procesa la información y genera el reporte solicitado. \\
\hline
Postcondiciones & El reporte queda disponible para su consulta o descarga. \\
\hline
\end{tabular}
\end{center}
\end{minipage}

\begin{minipage}{\linewidth}
\bigskip
\noindent\textbf{RF-12. Identificación de productividad por parcela}
\label{tab:RF-12}

\begin{center}
\begin{tabular}{|p{3.5cm}|p{10.7cm}|}
\hline
Identificador & RF-12 \\
\hline
Nombre & Identificación de productividad por parcela \\
\hline
Descripción & El sistema deberá permitir identificar las parcelas con mayor y menor productividad mediante el análisis de la producción registrada. \\
\hline
Actor / origen & Administrador: EV-01 [07:54], [08:02]. \\
\hline
Entradas & Período de análisis (fecha, obligatorio); vista derivada del resultado de la consulta de producción por período (RF-07). \\
\hline
Salidas & Listado de parcelas ordenado por producción, identificando la de mayor y la de menor productividad. \\
\hline
Prioridad (MoSCoW) & Debería tener. \\
\hline
Criterio de verificación & El sistema identifica correctamente la parcela de mayor y la de menor producción del período en el 100~\% de los casos, verificado sobre un conjunto de prueba de 20 parcelas con producción conocida. \\
\hline
Precondiciones & Deben existir registros históricos de producción por parcela. \\
\hline
Flujo principal & 1. El administrador accede al módulo de análisis. 2. El sistema procesa la información histórica. 3. Se presentan los resultados de productividad por parcela. \\
\hline
Postcondiciones & El administrador dispone de información para evaluar el rendimiento de las parcelas. \\
\hline
\end{tabular}
\end{center}
\end{minipage}

\begin{minipage}{\linewidth}
\bigskip
\noindent\textbf{RF-13. Generación de estadísticas agrícolas}
\label{tab:RF-13}

\begin{center}
\begin{tabular}{|p{3.5cm}|p{10.7cm}|}
\hline
Identificador & RF-13 \\
\hline
Nombre & Generación de estadísticas agrícolas \\
\hline
Descripción & El sistema deberá generar estadísticas basadas en la información registrada, permitiendo visualizar indicadores relacionados con la producción, las actividades realizadas y el rendimiento de los cultivos. \\
\hline
Actor / origen & Administrador: EV-01 [07:28], [07:36], [07:43]. \\
\hline
Entradas & Ninguna directa; procesa la información ya registrada en el sistema. \\
\hline
Salidas & Gráficos e indicadores de producción, actividades y rendimiento. \\
\hline
Prioridad (MoSCoW) & Debería tener. \\
\hline
Criterio de verificación & El sistema presenta las estadísticas agrícolas en $\leq$ 3,0 s con la información de un período de 12 meses cargada, verificado en el percentil 95 de 20 ejecuciones. \\
\hline
Precondiciones & Debe existir información registrada en la base de datos. \\
\hline
Flujo principal & 1. El administrador accede al módulo de estadísticas. 2. El sistema procesa los datos disponibles. 3. Se muestran gráficos e indicadores estadísticos. \\
\hline
Postcondiciones & Las estadísticas quedan disponibles para apoyar la toma de decisiones. \\
\hline
\end{tabular}
\end{center}
\end{minipage}

\begin{minipage}{\linewidth}
\bigskip
\noindent\textbf{RF-14. Gestión de usuarios}
\label{tab:RF-14}

\begin{center}
\begin{tabular}{|p{3.5cm}|p{10.7cm}|}
\hline
Identificador & RF-14 \\
\hline
Nombre & Gestión de usuarios \\
\hline
Descripción & El sistema deberá permitir registrar, modificar, consultar y desactivar las cuentas de los usuarios autorizados, asignando el rol correspondiente a cada uno. \\
\hline
Actor / origen & Administrador: R-02 (el acceso a la información dependerá del rol asignado a cada usuario); EV-01 [02:35]. \\
\hline
Entradas & Nombre de usuario (texto, obligatorio); rol (selección, obligatorio); estado (selección: activo o inactivo, obligatorio). \\
\hline
Salidas & Cuenta de usuario registrada, modificada o desactivada, con el rol correspondiente. \\
\hline
Prioridad (MoSCoW) & Debe tener. \\
\hline
Criterio de verificación & El 100~\% de los cambios de rol o estado de un usuario se reflejan en sus permisos de acceso en $\leq$ 2,0 s, verificado sobre una muestra de 20 cuentas de prueba. \\
\hline
Precondiciones & El administrador debe haber iniciado sesión. \\
\hline
Flujo principal & 1. El administrador accede al módulo de usuarios. 2. Registra o modifica la información del usuario. 3. El sistema valida los datos y guarda los cambios. \\
\hline
Postcondiciones & La información del usuario queda actualizada en el sistema. \\
\hline
\end{tabular}
\end{center}
\end{minipage}

\begin{minipage}{\linewidth}
\bigskip
\noindent\textbf{RF-15. Consultar recomendaciones de IA}
\label{tab:RF-15}

\begin{center}
\begin{tabular}{|p{3.5cm}|p{10.7cm}|}
\hline
Identificador & RF-15 \\
\hline
Nombre & Consultar recomendaciones de IA \\
\hline
Descripción & El sistema deberá analizar la información registrada sobre los cultivos de verde y cacao para generar recomendaciones que apoyen la planificación de actividades y la toma de decisiones por parte del administrador. \\
\hline
Actor / origen & Administrador: EV-01 [09:26], [09:36]. \\
\hline
Entradas & Ninguna directa; el módulo de IA procesa los registros históricos existentes de producción, actividades e inventario. \\
\hline
Salidas & Recomendaciones relacionadas con la gestión de los cultivos, sujetas a la restricción R-05. \\
\hline
Prioridad (MoSCoW) & Podría tener. \\
\hline
Criterio de verificación & El sistema presenta recomendaciones fundamentadas en la información registrada en $\leq$ 5,0 s, sin modificar automáticamente los datos almacenados, verificado sobre 20 solicitudes de recomendación. \\
\hline
Precondiciones & Deben existir registros históricos suficientes de producción, actividades e inventario. \\
\hline
Flujo principal & 1. El administrador accede al módulo de Inteligencia Artificial. 2. El sistema analiza la información disponible. 3. Se generan recomendaciones relacionadas con la gestión de los cultivos. \\
\hline
Postcondiciones & Las recomendaciones quedan disponibles para su consulta por parte del administrador. \\
\hline
\end{tabular}
\end{center}
\end{minipage}

\begin{minipage}{\linewidth}
\bigskip
\noindent\textbf{RF-16. Registro de enfermedades del cultivo}
\label{tab:RF-16}

\begin{center}
\begin{tabular}{|p{3.5cm}|p{10.7cm}|}
\hline
Identificador & RF-16 \\
\hline
Nombre & Registro de enfermedades del cultivo \\
\hline
Descripción & El sistema deberá permitir registrar las enfermedades detectadas en los cultivos de verde y cacao, indicando el cultivo afectado, el nombre de la enfermedad, el nivel de gravedad y observaciones, con el fin de llevar un control fitosanitario y facilitar la generación de alertas cuando sea necesario. \\
\hline
Actor / origen & Administrador, Trabajador agrícola: EV-02 [01:42]; EV-01 [01:16], [01:25]. \\
\hline
Entradas & Cultivo afectado (selección, obligatorio); nombre de la enfermedad (selección desde catálogo, RF-20, obligatorio); nivel de gravedad (selección, obligatorio); observaciones (texto, opcional). \\
\hline
Salidas & Enfermedad registrada y, si la gravedad es alta, alerta generada (RF-25). \\
\hline
Prioridad (MoSCoW) & Debería tener. \\
\hline
Criterio de verificación & El 100~\% de los registros de enfermedad con nivel de gravedad alto generan una alerta en $\leq$ 5,0 s desde su registro, verificado sobre una muestra de 30 casos. \\
\hline
Precondiciones & Debe existir un cultivo previamente registrado en el sistema. \\
\hline
Flujo principal & 1. El usuario selecciona el cultivo afectado. 2. Registra la enfermedad detectada. 3. Ingresa el nivel de gravedad y observaciones. 4. El sistema valida y almacena la información. 5. Si la gravedad es alta, el sistema genera una alerta. \\
\hline
Postcondiciones & La enfermedad queda registrada y disponible para su consulta. \\
\hline
\end{tabular}
\end{center}
\end{minipage}

A continuación se incorporan los once requisitos funcionales adicionales (RF-17 a RF-27), derivados de las capacidades identificadas en la evidencia de campo que no quedaban cubiertas por los dieciséis requisitos migrados. Cada uno sigue estrictamente la plantilla de ocho atributos, sin campos adicionales.

\begin{minipage}{\linewidth}
\bigskip
\noindent\textbf{RF-17. Registro de peso y comprobante de acopio}
\label{tab:RF-17}

\begin{center}
\begin{tabular}{|p{3.5cm}|p{10.7cm}|}
\hline
Identificador & RF-17 \\
\hline
Nombre & Registro de peso y comprobante de acopio \\
\hline
Descripción & El sistema deberá permitir registrar el peso entregado y el número de comprobante emitido por el centro de acopio para cada cosecha, como respaldo de la producción declarada. \\
\hline
Actor / origen & Trabajador agrícola, Administrador: EV-01 [04:32], [04:47]. \\
\hline
Entradas & Peso entregado (numérico, kg o toneladas, obligatorio); número de comprobante (texto, obligatorio). \\
\hline
Salidas & Registro de cosecha (RF-05) con peso y comprobante asociados. \\
\hline
Prioridad (MoSCoW) & Debería tener. \\
\hline
Criterio de verificación & El 100~\% de los registros de cosecha con comprobante de acopio permiten recuperar el peso y el número de comprobante en una consulta posterior, verificado sobre una muestra de 30 registros. \\
\hline
\end{tabular}
\end{center}
\end{minipage}

\begin{minipage}{\linewidth}
\bigskip
\noindent\textbf{RF-18. Comparación de producción entre periodos}
\label{tab:RF-18}

\begin{center}
\begin{tabular}{|p{3.5cm}|p{10.7cm}|}
\hline
Identificador & RF-18 \\
\hline
Nombre & Comparación de producción entre periodos \\
\hline
Descripción & El sistema deberá permitir comparar la producción registrada entre dos periodos seleccionados por el administrador, para un mismo cultivo o parcela. \\
\hline
Actor / origen & Administrador: EV-01 [04:47], [04:54]. \\
\hline
Entradas & Periodo 1, periodo 2, cultivo o parcela (selección, obligatorio); vista derivada del resultado de la consulta de producción por período (RF-07). \\
\hline
Salidas & Diferencia porcentual y absoluta de producción entre los dos periodos seleccionados. \\
\hline
Prioridad (MoSCoW) & Debería tener. \\
\hline
Criterio de verificación & La comparación entre dos periodos con 500 registros de producción se presenta en $\leq$ 3,0 s, verificado en el percentil 95 de 30 ejecuciones. \\
\hline
\end{tabular}
\end{center}
\end{minipage}

\begin{minipage}{\linewidth}
\bigskip
\noindent\textbf{RF-19. Reporte de pérdidas por plagas}
\label{tab:RF-19}

\begin{center}
\begin{tabular}{|p{3.5cm}|p{10.7cm}|}
\hline
Identificador & RF-19 \\
\hline
Nombre & Reporte de pérdidas por plagas \\
\hline
Descripción & El sistema deberá generar un reporte de pérdidas de producción estimadas asociadas a plagas y enfermedades, por cultivo y por parcela. \\
\hline
Actor / origen & Administrador: EV-01 [08:17], [08:22], [08:29] (reporte señalado como el más importante por el administrador). \\
\hline
Entradas & Periodo (fecha, obligatorio); cultivo o parcela (selección, opcional). \\
\hline
Salidas & Reporte con pérdida estimada por plaga o enfermedad, cultivo y parcela. \\
\hline
Prioridad (MoSCoW) & Debe tener. \\
\hline
Criterio de verificación & El reporte de pérdidas por plagas se genera en $\leq$ 3,0 s con 500 registros de enfermedad cargados, verificado en el percentil 95 de 30 ejecuciones. \\
\hline
\end{tabular}
\end{center}
\end{minipage}

\begin{minipage}{\linewidth}
\bigskip
\noindent\textbf{RF-20. Catálogo de plagas y enfermedades por cultivo}
\label{tab:RF-20}

\begin{center}
\begin{tabular}{|p{3.5cm}|p{10.7cm}|}
\hline
Identificador & RF-20 \\
\hline
Nombre & Catálogo de plagas y enfermedades por cultivo \\
\hline
Descripción & El sistema deberá mantener un catálogo de plagas y enfermedades asociadas a cada tipo de cultivo, con nombre normalizado, para estandarizar el registro de incidencias fitosanitarias. \\
\hline
Actor / origen & Administrador: EV-02 [01:42]; EV-01 [01:16], [01:25]. \\
\hline
Entradas & Nombre de la plaga o enfermedad (texto, obligatorio); cultivo asociado (selección, obligatorio). \\
\hline
Salidas & Lista de plagas y enfermedades disponible para selección en el registro de enfermedades (RF-16). \\
\hline
Prioridad (MoSCoW) & Debería tener. \\
\hline
Criterio de verificación & El 100~\% de los registros de enfermedad (RF-16) permiten seleccionar el nombre desde el catálogo sin ingreso de texto libre, verificado sobre una muestra de 30 registros. \\
\hline
\end{tabular}
\end{center}
\end{minipage}

\begin{minipage}{\linewidth}
\bigskip
\noindent\textbf{RF-21. Planificación de compra de insumos por calendario}
\label{tab:RF-21}

\begin{center}
\begin{tabular}{|p{3.5cm}|p{10.7cm}|}
\hline
Identificador & RF-21 \\
\hline
Nombre & Planificación de compra de insumos por calendario \\
\hline
Descripción & El sistema deberá permitir planificar la compra de insumos agrícolas asignando cantidades estimadas por cultivo a meses específicos del año. \\
\hline
Actor / origen & Administrador: EV-01 [05:36], [05:44]. \\
\hline
Entradas & Insumo (selección, obligatorio); cultivo (selección, obligatorio); mes (selección, obligatorio); cantidad estimada (numérico, obligatorio). \\
\hline
Salidas & Calendario de compras planificadas por insumo y mes. \\
\hline
Prioridad (MoSCoW) & Debería tener. \\
\hline
Criterio de verificación & El sistema muestra el calendario de compras planificadas de los 12 meses en $\leq$ 2,0 s, verificado en el percentil 95 de 20 ejecuciones. \\
\hline
\end{tabular}
\end{center}
\end{minipage}

\begin{minipage}{\linewidth}
\bigskip
\noindent\textbf{RF-22. Registro de entrega de insumos a trabajadores}
\label{tab:RF-22}

\begin{center}
\begin{tabular}{|p{3.5cm}|p{10.7cm}|}
\hline
Identificador & RF-22 \\
\hline
Nombre & Registro de entrega de insumos a trabajadores \\
\hline
Descripción & El sistema deberá permitir registrar la entrega de insumos agrícolas del administrador a un trabajador, actualizando la disponibilidad del inventario. \\
\hline
Actor / origen & Administrador: EV-02 [02:25]. \\
\hline
Entradas & Insumo (selección, obligatorio); cantidad entregada (numérico, obligatorio); trabajador receptor (selección, obligatorio); fecha (fecha, obligatorio). \\
\hline
Salidas & Actualización del inventario (RF-06) y registro de entrega asociado al trabajador. \\
\hline
Prioridad (MoSCoW) & Podría tener. \\
\hline
Criterio de verificación & El 100~\% de las entregas registradas descuentan correctamente la cantidad del inventario disponible, verificado sobre una muestra de 30 entregas. \\
\hline
\end{tabular}
\end{center}
\end{minipage}

\begin{minipage}{\linewidth}
\bigskip
\noindent\textbf{RF-23. Registro de ingresos}
\label{tab:RF-23}

\begin{center}
\begin{tabular}{|p{3.5cm}|p{10.7cm}|}
\hline
Identificador & RF-23 \\
\hline
Nombre & Registro de ingresos \\
\hline
Descripción & El sistema deberá permitir registrar los ingresos generados por la venta de la producción agrícola, asociados a cultivo y periodo. \\
\hline
Actor / origen & Administrador: EV-01 [06:23], [06:28], [06:37]. \\
\hline
Entradas & Monto (numérico, obligatorio); cultivo (selección, obligatorio); periodo (fecha, obligatorio); fecha (fecha, obligatorio). \\
\hline
Salidas & Registro de ingreso disponible para consulta por periodo. \\
\hline
Prioridad (MoSCoW) & Podría tener (necesidad secundaria frente a las prioridades explícitas del administrador). \\
\hline
Criterio de verificación & El registro de un ingreso queda disponible para consulta en $\leq$ 2,0 s tras su ingreso, verificado sobre 20 registros consecutivos. \\
\hline
\end{tabular}
\end{center}
\end{minipage}

\begin{minipage}{\linewidth}
\bigskip
\noindent\textbf{RF-24. Consulta de historial de registros por trabajador y fecha}
\label{tab:RF-24}

\begin{center}
\begin{tabular}{|p{3.5cm}|p{10.7cm}|}
\hline
Identificador & RF-24 \\
\hline
Nombre & Consulta de historial de registros por trabajador y fecha \\
\hline
Descripción & El sistema deberá permitir consultar el historial de labores y cosechas registradas, filtrado por trabajador y por rango de fechas, para permitir la reconciliación de la información y evitar la desviación de datos identificada en el registro en papel. \\
\hline
Actor / origen & Administrador: EV-01 [02:04], [02:17]. \\
\hline
Entradas & Trabajador (selección, opcional); rango de fechas (obligatorio). \\
\hline
Salidas & Listado de labores y cosechas registradas en el rango filtrado. \\
\hline
Prioridad (MoSCoW) & Debe tener. \\
\hline
Criterio de verificación & La consulta del historial filtrado por trabajador y fecha con 500 registros se presenta en $\leq$ 2,0 s, verificado en el percentil 95 de 30 ejecuciones. \\
\hline
\end{tabular}
\end{center}
\end{minipage}

\begin{minipage}{\linewidth}
\bigskip
\noindent\textbf{RF-25. Generación de alertas automáticas de atención inmediata}
\label{tab:RF-25}

\begin{center}
\begin{tabular}{|p{3.5cm}|p{10.7cm}|}
\hline
Identificador & RF-25 \\
\hline
Nombre & Generación de alertas automáticas de atención inmediata \\
\hline
Descripción & El sistema deberá generar alertas automáticas visibles al administrador cuando se registre una condición que requiera atención inmediata, incluyendo enfermedades de gravedad alta e insumos con disponibilidad crítica. \\
\hline
Actor / origen & Administrador: EV-01 [10:33]. \\
\hline
Entradas & Ninguna directa; se genera a partir de los registros de RF-16 y RF-06. \\
\hline
Salidas & Alerta con tipo, origen (parcela o insumo) y momento de generación. \\
\hline
Prioridad (MoSCoW) & Debe tener. \\
\hline
Criterio de verificación & El 100~\% de las condiciones que requieren atención inmediata (enfermedades con gravedad alta registradas vía RF-16 e insumos con disponibilidad crítica registrados vía RF-06) generan una alerta visible en $\leq$ 5,0 s desde su registro, verificada cada fuente por separado sobre una muestra de 30 casos por tipo. \\
\hline
\end{tabular}
\end{center}
\end{minipage}

\begin{minipage}{\linewidth}
\bigskip
\noindent\textbf{RF-26. Detección temprana de señales de enfermedad en parcelas mediante IA}
\label{tab:RF-26}

\begin{center}
\begin{tabular}{|p{3.5cm}|p{10.7cm}|}
\hline
Identificador & RF-26 \\
\hline
Nombre & Detección temprana de señales de enfermedad en parcelas mediante IA \\
\hline
Descripción & El sistema deberá analizar mediante IA la información histórica registrada de una parcela para generar una alerta temprana cuando se identifiquen patrones asociados a una posible enfermedad, antes de su confirmación visual por el usuario. \\
\hline
Actor / origen & Administrador: EV-01 [09:02]. \\
\hline
Entradas & Registros históricos de la parcela (labores, cosechas, enfermedades previas). \\
\hline
Salidas & Alerta temprana de posible enfermedad, con parcela y nivel de confianza del modelo. \\
\hline
Prioridad (MoSCoW) & Debe tener. \\
\hline
Criterio de verificación & El módulo de IA genera la alerta temprana con una exactitud (accuracy) mínima del 70~\% sobre un conjunto de validación de al menos 30 casos etiquetados. \\
\hline
\end{tabular}
\end{center}
\end{minipage}

\begin{minipage}{\linewidth}
\bigskip
\noindent\textbf{RF-27. Consulta de producción agregada por cultivo}
\label{tab:RF-27}

\begin{center}
\begin{tabular}{|p{3.5cm}|p{10.7cm}|}
\hline
Identificador & RF-27 \\
\hline
Nombre & Consulta de producción agregada por cultivo \\
\hline
Descripción & El sistema deberá permitir consultar la producción total registrada agrupada por cultivo, sin necesidad de discriminar por parcela individual. \\
\hline
Actor / origen & Administrador: EV-01 [09:57]. \\
\hline
Entradas & Cultivo (selección); periodo (opcional); vista derivada del resultado de la consulta de producción agregada del período (RF-07). \\
\hline
Salidas & Producción total agregada del cultivo seleccionado. \\
\hline
Prioridad (MoSCoW) & Debería tener. \\
\hline
Criterio de verificación & La consulta de producción agregada por cultivo con 500 registros se presenta en $\leq$ 2,0 s, verificado en el percentil 95 de 30 ejecuciones. \\
\hline
\end{tabular}
\end{center}
\end{minipage}

A continuación se incorporan los cuatro requisitos funcionales que especifican los dos componentes de inteligencia artificial del sistema (RF-28 a RF-31), exigidos por el Paso 3 de la PE5. RF-28 completa el componente IA-01, ya existente en la línea base como RF-15 y RF-26, con el mecanismo de retroalimentación humana que exige la restricción R-05. RF-29 a RF-31 especifican IA-02, el segundo componente de inteligencia artificial incorporado en esta práctica: un detector de plagas y enfermedades por imagen, justificado y descrito en la Sección~\ref{sec:componentes-ia}. Los cuatro siguen estrictamente la plantilla de ocho atributos, sin campos adicionales.

\begin{minipage}{\linewidth}
\bigskip
\noindent\textbf{RF-28. Confirmar o descartar una alerta o recomendación del módulo de IA}
\label{tab:RF-28}

\begin{center}
\begin{tabular}{|p{3.5cm}|p{10.7cm}|}
\hline
Identificador & RF-28 \\
\hline
Nombre & Confirmar o descartar una alerta o recomendación del módulo de IA \\
\hline
Descripción & El sistema deberá permitir que el administrador confirme o descarte cada alerta temprana (RF-26) o recomendación (RF-15) presentada, registrando la decisión como retroalimentación del componente IA-01. \\
\hline
Actor / origen & Administrador: EV-01 [09:02], [09:10] (necesidad de verificar la información antes de actuar); R-05. \\
\hline
Entradas & Alerta o recomendación presentada (RF-15, RF-26); decisión del administrador (confirmar o descartar). \\
\hline
Salidas & Registro de la decisión, disponible como retroalimentación para el plan de monitoreo del componente (Sección~\ref{sec:componentes-ia}); ningún cambio se aplica automáticamente sin esta confirmación. \\
\hline
Prioridad (MoSCoW) & Debería tener. \\
\hline
Criterio de verificación & El 100~\% de las alertas y recomendaciones generadas quedan con una decisión de confirmación o descarte registrada, verificado sobre una muestra de 30 alertas o recomendaciones; ninguna acción se ejecuta automáticamente sin esta decisión. \\
\hline
\end{tabular}
\end{center}
\end{minipage}

\begin{minipage}{\linewidth}
\bigskip
\noindent\textbf{RF-29. Clasificar una fotografía de una incidencia fitosanitaria}
\label{tab:RF-29}

\begin{center}
\begin{tabular}{|p{3.5cm}|p{10.7cm}|}
\hline
Identificador & RF-29 \\
\hline
Nombre & Clasificar una fotografía de una incidencia fitosanitaria \\
\hline
Descripción & El sistema deberá, a partir de una fotografía tomada durante el registro de una incidencia fitosanitaria (RF-16), proponer una clase del catálogo de plagas y enfermedades del cultivo correspondiente (RF-20), o «sin hallazgo» / «no concluyente», junto con un nivel de confianza. \\
\hline
Actor / origen & Administrador, Trabajador agrícola: EV-01 [01:16], [01:25], [08:17]--[08:29] (pérdidas por plagas como reporte prioritario); EV-02 [01:34] (reporte de incidencias fitosanitarias por el trabajador agrícola). \\
\hline
Entradas & Fotografía tomada en campo; cultivo y parcela seleccionados en el registro de la incidencia (RF-16). \\
\hline
Salidas & Clase propuesta del catálogo (RF-20), o «sin hallazgo» / «no concluyente»; nivel de confianza; región de la imagen que motivó la propuesta (RNF-26). \\
\hline
Prioridad (MoSCoW) & Debería tener. \\
\hline
Criterio de verificación & Sobre un conjunto de 30 fotografías de campo con clase confirmada, el sistema devuelve una propuesta de clase y nivel de confianza para el 100~\% de las capturas, sin modificar automáticamente el registro de enfermedad (R-05). \\
\hline
\end{tabular}
\end{center}
\end{minipage}

\begin{minipage}{\linewidth}
\bigskip
\noindent\textbf{RF-30. Registrar la clase confirmada en el historial de la parcela}
\label{tab:RF-30}

\begin{center}
\begin{tabular}{|p{3.5cm}|p{10.7cm}|}
\hline
Identificador & RF-30 \\
\hline
Nombre & Registrar la clase confirmada en el historial de la parcela \\
\hline
Descripción & El sistema deberá permitir que el usuario confirme, corrija o descarte la propuesta de RF-29 y, solo tras esa decisión, completar el registro de la incidencia (RF-16) con la clase resultante. \\
\hline
Actor / origen & Administrador, Trabajador agrícola: RF-16; R-05. \\
\hline
Entradas & Propuesta de clase de RF-29; decisión del usuario (confirmar, corregir o descartar). \\
\hline
Salidas & Registro de enfermedad (RF-16) completado con la clase resultante y el origen «detección por imagen»; si se descarta, no se crea ni modifica ningún registro a partir de la propuesta. \\
\hline
Prioridad (MoSCoW) & Debería tener. \\
\hline
Criterio de verificación & El 100~\% de los registros de enfermedad con origen «detección por imagen» tienen una decisión explícita del usuario asociada, verificado sobre una muestra de 20 casos; ninguna propuesta descartada genera o modifica un registro. \\
\hline
\end{tabular}
\end{center}
\end{minipage}

\begin{minipage}{\linewidth}
\bigskip
\noindent\textbf{RF-31. Generar alerta por incidencia acumulada de una plaga o enfermedad por parcela}
\label{tab:RF-31}

\begin{center}
\begin{tabular}{|p{3.5cm}|p{10.7cm}|}
\hline
Identificador & RF-31 \\
\hline
Nombre & Generar alerta por incidencia acumulada de una plaga o enfermedad por parcela \\
\hline
Descripción & El sistema deberá generar una alerta mediante el mecanismo ya existente (RF-25) cuando el número de detecciones confirmadas (RF-30) de una misma clase en una misma parcela alcance un umbral de incidencia configurado por el administrador. \\
\hline
Actor / origen & Administrador: EV-01 [08:17]--[08:29] (pérdidas por plagas como preocupación prioritaria), [10:33] (alertas de atención inmediata). \\
\hline
Entradas & Detecciones confirmadas (RF-30); umbral de incidencia configurado por el administrador (valor por defecto propuesto: 3 detecciones en 7 días). \\
\hline
Salidas & Alerta de atención inmediata (RF-25) con parcela, clase y detecciones que la originaron. \\
\hline
Prioridad (MoSCoW) & Debería tener. \\
\hline
Criterio de verificación & El 100~\% de los casos de prueba en que se alcanza el umbral configurado generan alerta en $\leq$ 5,0 s, y ningún caso por debajo del umbral la genera, verificado sobre 10 escenarios simulados. \\
\hline
\end{tabular}
\end{center}
\end{minipage}

\subsection{Requisitos no funcionales}
\label{sec:requisitos-no-funcionales}

% Subseccion 3.2.1: Criterio de cuantificacion ISO/IEC 25010:2023
\subsubsection{Criterio de cuantificación ISO/IEC 25010:2023}
\label{sec:criterio-cuantificacion-rnf}

La cuantificación de los requisitos no funcionales del SGA sigue un criterio único: todo requisito no funcional debe expresarse mediante una métrica, un valor objetivo y un método de medición verificable, evitando criterios procedimentales que solo describan una actividad de prueba sin un umbral de aceptación. Este criterio se aplica sobre las nueve características del modelo de calidad de producto de la norma ISO/IEC 25010:2023~\cite{iso25010-2023}, ampliadas para esta práctica con la equidad como concepto operativo adicional exigido a los componentes de inteligencia artificial (Sección~4.4 de la guía de la PE5): ausencia de diferencias injustificadas de desempeño entre grupos, especificada con una métrica de brecha, los grupos comparados y la brecha máxima tolerada. La Tabla~\ref{tab:cobertura-rnf} resume la cobertura mínima exigida por característica y la cobertura alcanzada en el catálogo de la Sección~\ref{sec:catalogo-rnf}.

\begin{center}
\begin{longtable}{|p{5cm}|p{2.2cm}|p{6cm}|}
\hline
\textbf{Característica} & \textbf{Mínimo} & \textbf{RNF incluidos} \\
\hline
\endfirsthead
\hline
\textbf{Característica} & \textbf{Mínimo} & \textbf{RNF incluidos} \\
\hline
\endhead
\hline
\endfoot
\hline
\endlastfoot
Adecuación funcional & 1 & RNF-07, RNF-16, RNF-21 \\
\hline
Eficiencia de desempeño & 2 & RNF-01, RNF-08, RNF-17, RNF-22 \\
\hline
Compatibilidad & 1 & RNF-05 \\
\hline
Interacción de usuario & 2 & RNF-04, RNF-12 \\
\hline
Fiabilidad & 2 & RNF-02, RNF-09, RNF-18, RNF-23 \\
\hline
Seguridad & 3 & RNF-03, RNF-10, RNF-11 \\
\hline
Mantenibilidad & 1 & RNF-06 \\
\hline
Portabilidad & 1 & RNF-13 \\
\hline
Flexibilidad & 1 & RNF-14 \\
\hline
Explicabilidad (Sección~\ref{sec:explicabilidad-ia} y~\ref{sec:componentes-ia}) & 2 & RNF-15, RNF-26 \\
\hline
Equidad del componente de IA, Secc.~4.4 de la guía (Sección~\ref{sec:componentes-ia}) & 4 & RNF-19, RNF-20, RNF-24, RNF-25 \\
\hline
\end{longtable}
\label{tab:cobertura-rnf}
\end{center}

% Subseccion 3.2.2: Catalogo de requisitos no funcionales
\subsubsection{Catálogo de requisitos no funcionales}
\label{sec:catalogo-rnf}

Cada requisito no funcional se documenta mediante la ficha de diez atributos descrita en la Tabla~\ref{tab:plantilla-rnf}.

\begin{center}
\begin{longtable}{|p{3.8cm}|p{10.4cm}|}
\hline
\textbf{Atributo} & \textbf{Contenido} \\
\hline
\endfirsthead
\hline
\textbf{Atributo} & \textbf{Contenido} \\
\hline
\endhead
\hline
\endfoot
\hline
\endlastfoot
Identificador & Código único con el formato RNF-XX. \\
\hline
Nombre & Frase nominal breve que identifica la propiedad de calidad especificada. \\
\hline
Característica ISO/IEC 25010:2023 & Una de las nueve características del modelo de calidad de producto. \\
\hline
Subcaracterística & Subcaracterística específica dentro de la característica declarada. \\
\hline
Descripción & Enunciado en la forma «El sistema deberá\ldots», referido a una propiedad de calidad medible. \\
\hline
Métrica & Variable que se mide para verificar el cumplimiento del requisito. \\
\hline
Valor objetivo & Umbral numérico o condición que la métrica debe alcanzar. \\
\hline
Método de medición & Procedimiento, tamaño de muestra y condiciones bajo las cuales se mide la métrica. \\
\hline
Origen & Identificador de evidencia EV-XX con marca de tiempo, o referencia normativa cuando el requisito no se origina en una cita directa. \\
\hline
Prioridad & Alta, Media o Baja. \\
\hline
\end{longtable}
\label{tab:plantilla-rnf}
\end{center}

Los requisitos RNF-01 a RNF-06 corresponden a los identificados en la Entrega 2 (1B), migrados a esta ficha y corregidos: sus criterios de verificación originales describían una actividad de prueba sin un umbral cuantificable, y se sustituyen aquí por una métrica con valor objetivo y método de medición explícitos. La categoría «Usabilidad» de la Entrega 2 (1B) se renombra a «Interacción de usuario», conforme a la nomenclatura de la edición 2023 de la norma. Los requisitos RNF-07 a RNF-14 se incorporan para cubrir las características y los mínimos exigidos en la Entrega 3 (2A). El requisito RNF-15 se desarrolla íntegramente en la Sección~\ref{sec:explicabilidad-ia}. Los requisitos RNF-16 a RNF-26 se incorporan en la práctica PE5: RNF-16 a RNF-20 formalizan el rendimiento y la equidad del componente IA-01, ya existente en la línea base, y RNF-21 a RNF-26 especifican el rendimiento, la equidad y la explicabilidad del componente IA-02, incorporado en esta práctica (Sección~\ref{sec:componentes-ia}).

\begin{minipage}{\linewidth}
\bigskip
\noindent\textbf{RNF-01. Tiempo de respuesta en la consulta de parcelas}
\label{tab:RNF-01}

\begin{center}
\begin{tabular}{|p{4.2cm}|p{10cm}|}
\hline
Identificador & RNF-01 \\
\hline
Nombre & Tiempo de respuesta en la consulta de parcelas \\
\hline
Característica ISO/IEC 25010:2023 & Eficiencia de desempeño \\
\hline
Subcaracterística & Comportamiento temporal \\
\hline
Descripción & El sistema deberá responder a las consultas y operaciones del usuario dentro de un tiempo máximo bajo condiciones normales de uso. \\
\hline
Métrica & Tiempo de respuesta \\
\hline
Valor objetivo & El listado de parcelas responde en $\leq$ 2,0 s con 500 registros cargados. \\
\hline
Método de medición & Medición del percentil 95 de 50 ejecuciones. \\
\hline
Origen & ISO/IEC 25010:2023 (característica sin evidencia directa de campo); RF-02. \\
\hline
Prioridad & Alta \\
\hline
\end{tabular}
\end{center}
\end{minipage}

\begin{minipage}{\linewidth}
\bigskip
\noindent\textbf{RNF-02. Disponibilidad del sistema}
\label{tab:RNF-02}

\begin{center}
\begin{tabular}{|p{4.2cm}|p{10cm}|}
\hline
Identificador & RNF-02 \\
\hline
Nombre & Disponibilidad del sistema \\
\hline
Característica ISO/IEC 25010:2023 & Fiabilidad \\
\hline
Subcaracterística & Disponibilidad \\
\hline
Descripción & El sistema deberá estar disponible durante el horario de operación de la finca. \\
\hline
Métrica & Porcentaje de tiempo disponible (uptime) \\
\hline
Valor objetivo & Disponibilidad $\geq$ 99,0~\% mensual en horario 06:00--18:00. \\
\hline
Método de medición & Medida como uptime registrado durante 30 días. \\
\hline
Origen & ISO/IEC 25010:2023 (característica sin evidencia directa de campo). \\
\hline
Prioridad & Alta \\
\hline
\end{tabular}
\end{center}
\end{minipage}

\begin{minipage}{\linewidth}
\bigskip
\noindent\textbf{RNF-03. Control de acceso por rol}
\label{tab:RNF-03}

\begin{center}
\begin{tabular}{|p{4.2cm}|p{10cm}|}
\hline
Identificador & RNF-03 \\
\hline
Nombre & Control de acceso por rol \\
\hline
Característica ISO/IEC 25010:2023 & Seguridad \\
\hline
Subcaracterística & Autenticidad \\
\hline
Descripción & El sistema deberá permitir el acceso únicamente a usuarios autenticados mediante credenciales válidas y controlar los permisos según el rol asignado. \\
\hline
Métrica & Porcentaje de intentos de acceso fuera de rol rechazados \\
\hline
Valor objetivo & El 100~\% de los intentos de acceso a funciones fuera del rol asignado son rechazados. \\
\hline
Método de medición & Verificado sobre una matriz de 12 casos rol-función. \\
\hline
Origen & R-02, RD-04. \\
\hline
Prioridad & Alta \\
\hline
\end{tabular}
\end{center}
\end{minipage}

\begin{minipage}{\linewidth}
\bigskip
\noindent\textbf{RNF-04. Curva de aprendizaje de usuarios nuevos}
\label{tab:RNF-04}

\begin{center}
\begin{tabular}{|p{4.2cm}|p{10cm}|}
\hline
Identificador & RNF-04 \\
\hline
Nombre & Curva de aprendizaje de usuarios nuevos \\
\hline
Característica ISO/IEC 25010:2023 & Interacción de usuario \\
\hline
Subcaracterística & Capacidad de aprendizaje \\
\hline
Descripción & Las interfaces del sistema deberán ser suficientemente intuitivas para que un usuario nuevo realice las funciones básicas tras una capacitación breve. \\
\hline
Métrica & Porcentaje de usuarios de prueba que completan las funciones básicas tras la capacitación \\
\hline
Valor objetivo & El 100~\% de los usuarios de prueba completan las funciones básicas tras una capacitación de 30 minutos. \\
\hline
Método de medición & Prueba con al menos 5 usuarios nuevos, midiendo el tiempo necesario para ejecutar las tareas principales. \\
\hline
Origen & ISO/IEC 25010:2023, motivado por el perfil de usuario sin experiencia digital previa declarado en EV-02 [00:24]. \\
\hline
Prioridad & Media \\
\hline
\end{tabular}
\end{center}
\end{minipage}

\begin{minipage}{\linewidth}
\bigskip
\noindent\textbf{RNF-05. Compatibilidad con navegadores}
\label{tab:RNF-05}

\begin{center}
\begin{tabular}{|p{4.2cm}|p{10cm}|}
\hline
Identificador & RNF-05 \\
\hline
Nombre & Compatibilidad con navegadores \\
\hline
Característica ISO/IEC 25010:2023 & Compatibilidad \\
\hline
Subcaracterística & Coexistencia \\
\hline
Descripción & El sistema deberá funcionar correctamente en los navegadores Google Chrome, Microsoft Edge y Mozilla Firefox en sus versiones actuales. \\
\hline
Métrica & Proporción de casos de uso ejecutados sin error funcional, por navegador \\
\hline
Valor objetivo & 18 de 18 ejecuciones exitosas (6 casos de uso $\times$ 3 navegadores), 0 errores funcionales. \\
\hline
Método de medición & Ejecución de los seis casos de uso principales (UC-01 a UC-06) en cada uno de los tres navegadores compatibles. \\
\hline
Origen & RD-03. \\
\hline
Prioridad & Media \\
\hline
\end{tabular}
\end{center}
\end{minipage}

\begin{minipage}{\linewidth}
\bigskip
\noindent\textbf{RNF-06. Modularidad de la arquitectura}
\label{tab:RNF-06}

\begin{center}
\begin{tabular}{|p{4.2cm}|p{10cm}|}
\hline
Identificador & RNF-06 \\
\hline
Nombre & Modularidad de la arquitectura \\
\hline
Característica ISO/IEC 25010:2023 & Mantenibilidad \\
\hline
Subcaracterística & Modularidad \\
\hline
Descripción & El sistema deberá estar desarrollado utilizando una arquitectura modular que facilite la incorporación de nuevas funcionalidades y la corrección de errores. \\
\hline
Métrica & Número de módulos modificados al incorporar una funcionalidad de complejidad equivalente a un módulo existente \\
\hline
Valor objetivo & Una nueva funcionalidad (por ejemplo, un nuevo tipo de reporte) se incorpora modificando como máximo 3 módulos del sistema. \\
\hline
Método de medición & Revisión de la arquitectura documentada por un desarrollador independiente al cambio. \\
\hline
Origen & RD-05. \\
\hline
Prioridad & Media \\
\hline
\end{tabular}
\end{center}
\end{minipage}

\begin{minipage}{\linewidth}
\bigskip
\noindent\textbf{RNF-07. Cobertura de unidades de medida por cultivo}
\label{tab:RNF-07}

\begin{center}
\begin{tabular}{|p{4.2cm}|p{10cm}|}
\hline
Identificador & RNF-07 \\
\hline
Nombre & Cobertura de unidades de medida por cultivo \\
\hline
Característica ISO/IEC 25010:2023 & Adecuación funcional \\
\hline
Subcaracterística & Completitud funcional \\
\hline
Descripción & El sistema deberá permitir registrar la cosecha en la unidad de medida correspondiente a cada cultivo (racimos para plátano, tachos o quintales para cacao), sin forzar una unidad única para todos los cultivos. \\
\hline
Métrica & Porcentaje de registros de cosecha con la unidad de medida correcta según el cultivo \\
\hline
Valor objetivo & 100~\%. \\
\hline
Método de medición & Verificación sobre una muestra de 30 registros de cosecha de al menos dos cultivos distintos. \\
\hline
Origen & EV-02 [02:08]; EV-01 [04:47]. \\
\hline
Prioridad & Alta \\
\hline
\end{tabular}
\end{center}
\end{minipage}

\begin{minipage}{\linewidth}
\bigskip
\noindent\textbf{RNF-08. Tiempo de generación de reportes de pérdidas}
\label{tab:RNF-08}

\begin{center}
\begin{tabular}{|p{4.2cm}|p{10cm}|}
\hline
Identificador & RNF-08 \\
\hline
Nombre & Tiempo de generación de reportes de pérdidas \\
\hline
Característica ISO/IEC 25010:2023 & Eficiencia de desempeño \\
\hline
Subcaracterística & Comportamiento temporal \\
\hline
Descripción & El sistema deberá generar el reporte de pérdidas por plagas (RF-19) dentro de un tiempo máximo, con un volumen representativo de registros cargados. \\
\hline
Métrica & Tiempo de respuesta \\
\hline
Valor objetivo & $\leq$ 3,0 s con 500 registros de enfermedad cargados. \\
\hline
Método de medición & Medición del percentil 95 de 30 ejecuciones. \\
\hline
Origen & ISO/IEC 25010:2023 (característica sin evidencia directa de campo). \\
\hline
Prioridad & Alta \\
\hline
\end{tabular}
\end{center}
\end{minipage}

\begin{minipage}{\linewidth}
\bigskip
\noindent\textbf{RNF-09. Tolerancia a interrupciones durante el registro}
\label{tab:RNF-09}

\begin{center}
\begin{tabular}{|p{4.2cm}|p{10cm}|}
\hline
Identificador & RNF-09 \\
\hline
Nombre & Tolerancia a interrupciones durante el registro \\
\hline
Característica ISO/IEC 25010:2023 & Fiabilidad \\
\hline
Subcaracterística & Tolerancia a fallos \\
\hline
Descripción & El sistema deberá conservar sin pérdida los datos ingresados en un formulario de registro (cosecha, labor o enfermedad) ante una interrupción de conexión durante el envío, permitiendo su reintento. \\
\hline
Métrica & Porcentaje de registros recuperados tras una interrupción simulada de conexión \\
\hline
Valor objetivo & 100~\%. \\
\hline
Método de medición & Prueba de interrupción de conexión simulada sobre 20 intentos de registro. \\
\hline
Origen & ISO/IEC 25010:2023, motivado por la pérdida de datos declarada en EV-01 [02:04], [02:17]. \\
\hline
Prioridad & Alta \\
\hline
\end{tabular}
\end{center}
\end{minipage}

\begin{minipage}{\linewidth}
\bigskip
\noindent\textbf{RNF-10. Bloqueo por intentos fallidos de acceso}
\label{tab:RNF-10}

\begin{center}
\begin{tabular}{|p{4.2cm}|p{10cm}|}
\hline
Identificador & RNF-10 \\
\hline
Nombre & Bloqueo por intentos fallidos de acceso \\
\hline
Característica ISO/IEC 25010:2023 & Seguridad \\
\hline
Subcaracterística & Resistencia \\
\hline
Descripción & El sistema deberá bloquear temporalmente una cuenta de usuario tras un número determinado de intentos fallidos consecutivos de inicio de sesión. \\
\hline
Métrica & Número de intentos fallidos permitidos antes del bloqueo, y duración del bloqueo \\
\hline
Valor objetivo & 5 intentos; bloqueo mínimo de 15 minutos. \\
\hline
Método de medición & Prueba de 5 intentos fallidos consecutivos sobre 10 cuentas de prueba. \\
\hline
Origen & UC-01; RD-04. \\
\hline
Prioridad & Alta \\
\hline
\end{tabular}
\end{center}
\end{minipage}

\begin{minipage}{\linewidth}
\bigskip
\noindent\textbf{RNF-11. Cifrado de credenciales almacenadas}
\label{tab:RNF-11}

\begin{center}
\begin{tabular}{|p{4.2cm}|p{10cm}|}
\hline
Identificador & RNF-11 \\
\hline
Nombre & Cifrado de credenciales almacenadas \\
\hline
Característica ISO/IEC 25010:2023 & Seguridad \\
\hline
Subcaracterística & Confidencialidad \\
\hline
Descripción & El sistema deberá almacenar las contraseñas de los usuarios mediante una función de cifrado unidireccional, sin conservar la contraseña en texto plano en ningún medio de almacenamiento. \\
\hline
Métrica & Porcentaje de contraseñas almacenadas en formato cifrado \\
\hline
Valor objetivo & 100~\%. \\
\hline
Método de medición & Revisión técnica del esquema de almacenamiento de credenciales en la base de datos. \\
\hline
Origen & ISO/IEC 25010:2023; R-04. \\
\hline
Prioridad & Alta \\
\hline
\end{tabular}
\end{center}
\end{minipage}

\begin{minipage}{\linewidth}
\bigskip
\noindent\textbf{RNF-12. Accesibilidad para usuarios sin experiencia digital previa}
\label{tab:RNF-12}

\begin{center}
\begin{tabular}{|p{4.2cm}|p{10cm}|}
\hline
Identificador & RNF-12 \\
\hline
Nombre & Accesibilidad para usuarios sin experiencia digital previa \\
\hline
Característica ISO/IEC 25010:2023 & Interacción de usuario \\
\hline
Subcaracterística & Inclusividad \\
\hline
Descripción & El sistema deberá permitir que un usuario sin experiencia previa en aplicaciones móviles complete el registro de una cosecha sin asistencia de un tercero. \\
\hline
Métrica & Porcentaje de usuarios de prueba sin experiencia previa que completan el registro sin asistencia \\
\hline
Valor objetivo & $\geq$ 80~\%. \\
\hline
Método de medición & Prueba de usabilidad con al menos 5 usuarios de prueba sin experiencia previa en aplicaciones. \\
\hline
Origen & EV-02 [00:24], [02:33], [02:43]. \\
\hline
Prioridad & Media \\
\hline
\end{tabular}
\end{center}
\end{minipage}

\begin{minipage}{\linewidth}
\bigskip
\noindent\textbf{RNF-13. Acceso desde dispositivo móvil}
\label{tab:RNF-13}

\begin{center}
\begin{tabular}{|p{4.2cm}|p{10cm}|}
\hline
Identificador & RNF-13 \\
\hline
Nombre & Acceso desde dispositivo móvil \\
\hline
Característica ISO/IEC 25010:2023 & Portabilidad \\
\hline
Subcaracterística & Adaptabilidad \\
\hline
Descripción & El sistema deberá ser accesible y operable desde un navegador móvil en un teléfono inteligente, sin requerir instalación de una aplicación nativa. \\
\hline
Métrica & Porcentaje de funciones principales operables desde un navegador móvil \\
\hline
Valor objetivo & 100~\% de las funciones de RF-01, RF-04 y RF-05. \\
\hline
Método de medición & Prueba funcional sobre un navegador móvil en al menos 2 tamaños de pantalla. \\
\hline
Origen & EV-01 [08:42]. \\
\hline
Prioridad & Alta \\
\hline
\end{tabular}
\end{center}
\end{minipage}

\begin{minipage}{\linewidth}
\bigskip
\noindent\textbf{RNF-14. Incorporación de nuevos cultivos sin modificar código}
\label{tab:RNF-14}

\begin{center}
\begin{tabular}{|p{4.2cm}|p{10cm}|}
\hline
Identificador & RNF-14 \\
\hline
Nombre & Incorporación de nuevos cultivos sin modificar código \\
\hline
Característica ISO/IEC 25010:2023 & Flexibilidad \\
\hline
Subcaracterística & Adaptabilidad \\
\hline
Descripción & El sistema deberá permitir agregar un nuevo tipo de cultivo y su unidad de medida asociada sin requerir modificaciones al código fuente de la aplicación. \\
\hline
Métrica & Tiempo requerido para incorporar un nuevo tipo de cultivo mediante configuración \\
\hline
Valor objetivo & $\leq$ 30 minutos, sin intervención de un desarrollador. \\
\hline
Método de medición & Prueba de incorporación de un cultivo de prueba por parte de un administrador capacitado. \\
\hline
Origen & ISO/IEC 25010:2023; RD-05. \\
\hline
Prioridad & Media \\
\hline
\end{tabular}
\end{center}
\end{minipage}

\bigskip
\noindent RNF-15, correspondiente a la explicabilidad de las alertas y recomendaciones del módulo de IA, se documenta íntegramente en la Sección~\ref{sec:explicabilidad-ia}, ficha de la Tabla~\ref{tab:RNF-15}.

\subsection{Requisitos no funcionales de los componentes de inteligencia artificial}
\label{sec:explicabilidad-ia}

Esta sección desarrolla, con el nivel de detalle que exige la Sección 4.4 de la guía de la PE5, los requisitos no funcionales específicos de los dos componentes de inteligencia artificial del sistema: las decisiones automáticas que producen (3.3.1), las necesidades de explicación de cada tipo de usuario (3.3.2) y la ficha completa de RNF-15 a RNF-26 (3.3.3), que cubren rendimiento, equidad y explicabilidad. El origen histórico de esta sección es el requisito de explicabilidad de IA-01 (RNF-15); por eso se desarrolla primero en detalle el marco conceptual de explicabilidad, aplicable también a RNF-26 (IA-02).

El módulo de Inteligencia Artificial del SGA genera alertas y recomendaciones que inciden en decisiones sobre el manejo de los cultivos. Por tratarse de un componente de IA, el sistema incorpora un requisito de explicabilidad, entendida como la capacidad del sistema de proporcionar información comprensible sobre cómo y por qué llegó a una salida determinada~\cite{chazette2020}. El marco de referencia utilizado es el modelo de explicabilidad de Chazette, Brunotte y Speith~\cite{chazette2021,chazette2022}.

La evidencia central proviene de EV-01 [09:02] y [09:10]: el administrador condiciona su confianza en el sistema a la posibilidad de verificar la información antes de actuar, y reconoce explícitamente que la Inteligencia Artificial tiene margen de error. De esta evidencia se derivan cuatro rasgos de la necesidad de explicación: el usuario exige poder \textbf{verificar} la salida antes de actuar; debe comunicarse la \textbf{incertidumbre} asociada, no solo el resultado; el tipo de explicación requerido es \textbf{causal y por evidencia}, referida a los datos en que se basó la alerta para poder contrastarla en campo; y el momento de la explicación debe ser \textbf{simultáneo a la alerta}, no posterior. El perfil del jornalero, que no utiliza aplicaciones (EV-02 [02:33]), exige explicaciones de menor densidad que las del administrador. Este requisito se enlaza con la restricción R-05 de la Entrega 2 (1B): las recomendaciones de la IA apoyan la decisión y no sustituyen el criterio del usuario.

% Subseccion 3.3.1: Decisiones automaticas del modulo de IA
\subsubsection{Decisiones automáticas del módulo de IA}
\label{sec:decisiones-ia}

\begin{center}
\begin{longtable}{|p{1.6cm}|p{5.5cm}|p{2.3cm}|p{4.3cm}|}
\hline
\textbf{ID} & \textbf{Decisión automática} & \textbf{RF asociado} & \textbf{Nivel de autonomía} \\
\hline
\endfirsthead
\hline
\textbf{ID} & \textbf{Decisión automática} & \textbf{RF asociado} & \textbf{Nivel de autonomía} \\
\hline
\endhead
\hline
\endfoot
\hline
\endlastfoot
D-01 & Generación de alerta temprana de posible enfermedad en una parcela & RF-26 & Sugerencia; requiere confirmación humana (R-05) \\
\hline
D-02 & Generación de recomendación de manejo del cultivo & RF-15 & Sugerencia; requiere confirmación humana (R-05) \\
\hline
D-03 & Cálculo del nivel de confianza asociado a una alerta o recomendación & RF-26, RF-15 & Automático; informativo, no accionable por sí solo \\
\hline
D-04 & Clasificación de una fotografía de campo en una clase de plaga o enfermedad del catálogo (o «sin hallazgo» / «no concluyente») & RF-29 & Sugerencia; el registro requiere confirmación humana (RF-30, R-05) \\
\hline
D-05 & Generación de alerta por incidencia acumulada de una plaga o enfermedad en una parcela & RF-31 & Automático sobre detecciones ya confirmadas; requiere atención humana, no ejecuta ninguna acción por sí sola (R-05) \\
\hline
\end{longtable}
\label{tab:decisiones-ia}
\end{center}

Las decisiones D-04 y D-05 corresponden al componente IA-02, incorporado en esta práctica (Sección~\ref{sec:componentes-ia}); ambas siguen el mismo principio que D-01 a D-03: la inteligencia artificial propone, el usuario decide.

% Subseccion 3.3.2: Necesidades de explicacion por tipo de usuario
\subsubsection{Necesidades de explicación por tipo de usuario}
\label{sec:necesidades-explicacion}

\begin{center}
\begin{longtable}{|p{2cm}|p{2.3cm}|p{5.3cm}|p{3.6cm}|}
\hline
\textbf{Decisión} & \textbf{Usuario afectado} & \textbf{Tipo de explicación requerida} & \textbf{Momento} \\
\hline
\endfirsthead
\hline
\textbf{Decisión} & \textbf{Usuario afectado} & \textbf{Tipo de explicación requerida} & \textbf{Momento} \\
\hline
\endhead
\hline
\endfoot
\hline
\endlastfoot
D-01 & Administrador & Causal y por evidencia: datos históricos y patrón que motivaron la alerta, con nivel de incertidumbre explícito & Simultáneo a la alerta (EV-01 [09:02]) \\
\hline
D-02 & Administrador & Causal y por evidencia, con nivel de incertidumbre explícito & Simultáneo a la recomendación (EV-01 [09:10]) \\
\hline
D-01, D-02 (tarea derivada) & Jornalero & Instructiva, de baja densidad: qué tarea ejecutar y en qué parcela, sin el detalle causal del modelo & Al recibir la tarea asignada, RF-09 (EV-02 [02:33]) \\
\hline
\end{longtable}
\label{tab:necesidades-explicacion}
\end{center}

% Subseccion 3.3.3: Fichas de los RNF de los componentes de IA
\subsubsection{Fichas de los requisitos no funcionales de rendimiento, equidad y explicabilidad (RNF-15 a RNF-26)}
\label{sec:ficha-explicabilidad}

\begin{minipage}{\linewidth}
\bigskip
\noindent\textbf{RNF-15. Explicabilidad de las alertas y recomendaciones del módulo de IA}
\label{tab:RNF-15}

\begin{center}
\begin{tabular}{|p{4.2cm}|p{10cm}|}
\hline
Identificador & RNF-15 \\
\hline
Nombre & Explicabilidad de las alertas y recomendaciones del módulo de IA \\
\hline
Característica ISO/IEC 25010:2023 & Interacción de usuario \\
\hline
Subcaracterística & Autodescriptividad \\
\hline
Descripción & El sistema deberá presentar, junto con cada alerta o recomendación generada por el módulo de IA, una explicación de los datos y el patrón que la originaron, junto con el nivel de incertidumbre asociado, para permitir que el administrador la verifique antes de actuar sobre ella. \\
\hline
Métrica & Cuatro métricas independientes: M15.1 proporción de alertas y recomendaciones que incluyen explicación causal y nivel de incertidumbre en su presentación; M15.2 tiempo adicional de presentación de la explicación; M15.3 extensión de la explicación; M15.4 comprensión reportada por el usuario. \\
\hline
Valor objetivo & M15.1 = 100~\% de las alertas y recomendaciones incluyen explicación causal y nivel de incertidumbre; M15.2 $\leq$ 1,0 s adicional respecto a la alerta; M15.3 extensión máxima de 3 líneas de texto; M15.4 comprensión reportada $\geq$ 4 sobre 5 en escala Likert. \\
\hline
Método de medición & Cada métrica se verifica por separado: M15.1 a M15.3 sobre una muestra de 20 alertas o recomendaciones generadas, y M15.4 mediante encuesta de comprensión en escala Likert de 5 puntos aplicada al administrador tras presentarle 10 alertas. \\
\hline
Origen & EV-01 [09:02], [09:10]; Chazette y Schneider (2020)~\cite{chazette2020}; Chazette, Brunotte y Speith (2021, 2022)~\cite{chazette2021,chazette2022}; R-05. \\
\hline
Prioridad & Alta \\
\hline
\end{tabular}
\end{center}
\end{minipage}

Los requisitos RNF-16 a RNF-26 formalizan, con ficha de diez atributos, el rendimiento y la equidad de los dos componentes de inteligencia artificial del sistema, y la explicabilidad de IA-02. Se documentan en detalle en la Sección~\ref{sec:componentes-ia}.

\begin{minipage}{\linewidth}
\bigskip
\noindent\textbf{RNF-16. Exactitud de la detección temprana de enfermedades (IA-01)}
\label{tab:RNF-16}

\begin{center}
\begin{tabular}{|p{4.2cm}|p{10cm}|}
\hline
Identificador & RNF-16 \\
\hline
Nombre & Exactitud de la detección temprana de enfermedades (IA-01) \\
\hline
Característica ISO/IEC 25010:2023 & Adecuación funcional \\
\hline
Subcaracterística & Corrección funcional \\
\hline
Descripción & El componente de detección temprana (RF-26) deberá alcanzar una exactitud mínima sobre un conjunto de validación etiquetado por confirmación en campo, antes de cada despliegue. \\
\hline
Métrica & Exactitud (\emph{accuracy}) \\
\hline
Valor objetivo & $\geq$ 70~\%. \\
\hline
Método de medición & Medido sobre un conjunto de validación de al menos 30 casos etiquetados por confirmación del administrador (RF-28), antes de cada despliegue. \\
\hline
Origen & RF-26; EV-01 [09:02]. \\
\hline
Prioridad & Alta \\
\hline
\end{tabular}
\end{center}
\end{minipage}

\begin{minipage}{\linewidth}
\bigskip
\noindent\textbf{RNF-17. Tiempo de respuesta del módulo de IA de alertas y recomendaciones (IA-01)}
\label{tab:RNF-17}

\begin{center}
\begin{tabular}{|p{4.2cm}|p{10cm}|}
\hline
Identificador & RNF-17 \\
\hline
Nombre & Tiempo de respuesta del módulo de IA de alertas y recomendaciones (IA-01) \\
\hline
Característica ISO/IEC 25010:2023 & Eficiencia de desempeño \\
\hline
Subcaracterística & Comportamiento temporal \\
\hline
Descripción & El sistema deberá presentar la alerta o recomendación generada por IA-01 dentro de un tiempo máximo desde la solicitud o el evento que la origina. \\
\hline
Métrica & Tiempo de respuesta \\
\hline
Valor objetivo & $\leq$ 5,0 s. \\
\hline
Método de medición & Percentil 95 de 30 solicitudes o eventos, con al menos 500 registros históricos cargados. \\
\hline
Origen & RF-15; RF-26. \\
\hline
Prioridad & Alta \\
\hline
\end{tabular}
\end{center}
\end{minipage}

\begin{minipage}{\linewidth}
\bigskip
\noindent\textbf{RNF-18. Disponibilidad del módulo de IA con degradación controlada (IA-01)}
\label{tab:RNF-18}

\begin{center}
\begin{tabular}{|p{4.2cm}|p{10cm}|}
\hline
Identificador & RNF-18 \\
\hline
Nombre & Disponibilidad del módulo de IA con degradación controlada (IA-01) \\
\hline
Característica ISO/IEC 25010:2023 & Fiabilidad \\
\hline
Subcaracterística & Disponibilidad \\
\hline
Descripción & Si el módulo de IA no está disponible o falla, el sistema deberá seguir operando sin recomendaciones ni alertas tempranas, sin bloquear ninguna función de registro o consulta. \\
\hline
Métrica & Porcentaje de funciones no relacionadas con IA operativas durante una interrupción simulada del módulo \\
\hline
Valor objetivo & 100~\%. \\
\hline
Método de medición & Prueba de interrupción simulada del módulo de IA sobre las funciones de registro (RF-04, RF-05, RF-16) y consulta (RF-07, RF-11), verificando que ninguna quede bloqueada. \\
\hline
Origen & ISO/IEC 25010:2023~\cite{iso25010-2023}; RNF-02. \\
\hline
Prioridad & Alta \\
\hline
\end{tabular}
\end{center}
\end{minipage}

\begin{minipage}{\linewidth}
\bigskip
\noindent\textbf{RNF-19. Equidad de desempeño entre cultivos de cacao y verde (IA-01)}
\label{tab:RNF-19}

\begin{center}
\begin{tabular}{|p{4.2cm}|p{10cm}|}
\hline
Identificador & RNF-19 \\
\hline
Nombre & Equidad de desempeño entre cultivos de cacao y verde (IA-01) \\
\hline
Característica ISO/IEC 25010:2023 & Equidad (concepto operativo de la PE5, Sección 4.4 de la guía; no es una de las nueve características de la norma) \\
\hline
Subcaracterística & Ausencia de diferencias injustificadas de desempeño entre grupos \\
\hline
Descripción & La exactitud de la detección temprana (RF-26) no deberá diferir de forma sustancial entre los registros de cultivo de cacao y los de plátano verde. \\
\hline
Métrica & Brecha de exactitud entre grupos (cacao frente a verde) \\
\hline
Valor objetivo & $\leq$ 10 puntos porcentuales. Tolerancia mayor que entre subgrupos de un mismo cultivo (RNF-20), porque cacao y verde son dominios agronómicos distintos con variabilidad propia esperable. \\
\hline
Método de medición & Exactitud de RF-26 calculada por separado sobre los registros de cada cultivo, con al menos 30 casos etiquetados por cultivo (RF-28); comparación trimestral. \\
\hline
Origen & RF-03 (el sistema cubre ambos cultivos); sesgo de densidad de datos entre cultivos documentado por el equipo (\texttt{fichas\_IA\_datos\_y\_monitoreo.md}). \\
\hline
Prioridad & Media \\
\hline
\end{tabular}
\end{center}
\end{minipage}

\begin{minipage}{\linewidth}
\bigskip
\noindent\textbf{RNF-20. Equidad de desempeño entre parcelas con distinto volumen de registro (IA-01)}
\label{tab:RNF-20}

\begin{center}
\begin{tabular}{|p{4.2cm}|p{10cm}|}
\hline
Identificador & RNF-20 \\
\hline
Nombre & Equidad de desempeño entre parcelas con distinto volumen de registro (IA-01) \\
\hline
Característica ISO/IEC 25010:2023 & Equidad (concepto operativo de la PE5, Sección 4.4 de la guía; no es una de las nueve características de la norma) \\
\hline
Subcaracterística & Ausencia de diferencias injustificadas de desempeño entre grupos \\
\hline
Descripción & La tasa de falsos positivos de las alertas tempranas (RF-26) no deberá diferir de forma sustancial entre parcelas con alto y bajo volumen histórico de registros. \\
\hline
Métrica & Brecha de tasa de falsos positivos entre parcelas de alto y bajo volumen de registro, usando la mediana de registros acumulados como punto de corte \\
\hline
Valor objetivo & $\leq$ 5 puntos porcentuales. \\
\hline
Método de medición & Cálculo trimestral sobre las alertas confirmadas o descartadas (RF-28), separando parcelas por encima y por debajo de la mediana de registros acumulados. \\
\hline
Origen & RF-15; RF-26; sesgo de subregistro documentado por el equipo: las parcelas revisadas con más frecuencia acumulan más incidencias (\texttt{fichas\_IA\_datos\_y\_monitoreo.md}). \\
\hline
Prioridad & Media \\
\hline
\end{tabular}
\end{center}
\end{minipage}

\begin{minipage}{\linewidth}
\bigskip
\noindent\textbf{RNF-21. Corrección de la clasificación de imágenes (IA-02)}
\label{tab:RNF-21}

\begin{center}
\begin{tabular}{|p{4.2cm}|p{10cm}|}
\hline
Identificador & RNF-21 \\
\hline
Nombre & Corrección de la clasificación de imágenes (IA-02) \\
\hline
Característica ISO/IEC 25010:2023 & Adecuación funcional \\
\hline
Subcaracterística & Corrección funcional \\
\hline
Descripción & El clasificador de fotografías (RF-29) deberá alcanzar un desempeño mínimo antes de cada despliegue, evaluado sobre un conjunto de prueba local no usado en el entrenamiento. \\
\hline
Métrica & F1 macro promediado por clase \\
\hline
Valor objetivo & $\geq$ 0,85; ninguna clase individual por debajo de 0,75. \\
\hline
Método de medición & Evaluado sobre un conjunto de prueba retenido de al menos 100 imágenes por clase, antes de cada despliegue. \\
\hline
Origen & RF-29; RF-20 (catálogo de clases). \\
\hline
Prioridad & Alta \\
\hline
\end{tabular}
\end{center}
\end{minipage}

\begin{minipage}{\linewidth}
\bigskip
\noindent\textbf{RNF-22. Tiempo de inferencia de la clasificación de imágenes (IA-02)}
\label{tab:RNF-22}

\begin{center}
\begin{tabular}{|p{4.2cm}|p{10cm}|}
\hline
Identificador & RNF-22 \\
\hline
Nombre & Tiempo de inferencia de la clasificación de imágenes (IA-02) \\
\hline
Característica ISO/IEC 25010:2023 & Eficiencia de desempeño \\
\hline
Subcaracterística & Comportamiento temporal \\
\hline
Descripción & El sistema deberá devolver la propuesta de clase de RF-29 dentro de un tiempo máximo desde que se toma o carga la fotografía. \\
\hline
Métrica & Tiempo de respuesta \\
\hline
Valor objetivo & $\leq$ 2,0 s. \\
\hline
Método de medición & Percentil 95 sobre 30 clasificaciones consecutivas. \\
\hline
Origen & RF-29. \\
\hline
Prioridad & Alta \\
\hline
\end{tabular}
\end{center}
\end{minipage}

\begin{minipage}{\linewidth}
\bigskip
\noindent\textbf{RNF-23. Continuidad del registro sin conectividad (IA-02)}
\label{tab:RNF-23}

\begin{center}
\begin{tabular}{|p{4.2cm}|p{10cm}|}
\hline
Identificador & RNF-23 \\
\hline
Nombre & Continuidad del registro sin conectividad (IA-02) \\
\hline
Característica ISO/IEC 25010:2023 & Fiabilidad \\
\hline
Subcaracterística & Tolerancia a fallos \\
\hline
Descripción & Si no hay conectividad o el módulo de clasificación no está disponible, el sistema deberá permitir completar el registro de la incidencia (RF-16) mediante selección manual del catálogo (RF-20), sin bloquear el registro. \\
\hline
Métrica & Porcentaje de registros de incidencia completables sin el módulo de clasificación disponible \\
\hline
Valor objetivo & 100~\%. \\
\hline
Método de medición & Prueba de indisponibilidad simulada del módulo sobre 20 intentos de registro de incidencia. \\
\hline
Origen & RF-16; RF-20; ISO/IEC 25010:2023~\cite{iso25010-2023}. \\
\hline
Prioridad & Alta \\
\hline
\end{tabular}
\end{center}
\end{minipage}

\begin{minipage}{\linewidth}
\bigskip
\noindent\textbf{RNF-24. Equidad de desempeño entre condiciones de luz (IA-02)}
\label{tab:RNF-24}

\begin{center}
\begin{tabular}{|p{4.2cm}|p{10cm}|}
\hline
Identificador & RNF-24 \\
\hline
Nombre & Equidad de desempeño entre condiciones de luz (IA-02) \\
\hline
Característica ISO/IEC 25010:2023 & Equidad (concepto operativo de la PE5, Sección 4.4 de la guía; no es una de las nueve características de la norma) \\
\hline
Subcaracterística & Ausencia de diferencias injustificadas de desempeño entre grupos \\
\hline
Descripción & El F1 del clasificador (RF-29) no deberá diferir de forma sustancial entre fotografías tomadas con luz plena y con sombra. \\
\hline
Métrica & Brecha de F1 entre condición de luz (plena frente a sombra) \\
\hline
Valor objetivo & $\leq$ 5 puntos porcentuales. \\
\hline
Método de medición & F1 calculado por separado sobre subconjuntos etiquetados por condición de luz dentro del conjunto de prueba de RNF-21, con al menos 50 imágenes por condición. \\
\hline
Origen & RF-29; sesgo de condiciones de captura documentado por el equipo (\texttt{fichas\_IA\_datos\_y\_monitoreo.md}). \\
\hline
Prioridad & Media \\
\hline
\end{tabular}
\end{center}
\end{minipage}

\begin{minipage}{\linewidth}
\bigskip
\noindent\textbf{RNF-25. Equidad de desempeño entre gamas de dispositivo (IA-02)}
\label{tab:RNF-25}

\begin{center}
\begin{tabular}{|p{4.2cm}|p{10cm}|}
\hline
Identificador & RNF-25 \\
\hline
Nombre & Equidad de desempeño entre gamas de dispositivo (IA-02) \\
\hline
Característica ISO/IEC 25010:2023 & Equidad (concepto operativo de la PE5, Sección 4.4 de la guía; no es una de las nueve características de la norma) \\
\hline
Subcaracterística & Ausencia de diferencias injustificadas de desempeño entre grupos \\
\hline
Descripción & El F1 del clasificador (RF-29) no deberá diferir de forma sustancial entre fotografías capturadas con dispositivos de gama alta y de gama baja. \\
\hline
Métrica & Brecha de F1 entre gama de dispositivo \\
\hline
Valor objetivo & $\leq$ 5 puntos porcentuales. \\
\hline
Método de medición & F1 calculado por separado sobre subconjuntos etiquetados por gama de dispositivo dentro del conjunto de prueba de RNF-21, con al menos 50 imágenes por gama. \\
\hline
Origen & RF-29; sesgo de predominio de un único dispositivo en la captura inicial, documentado por el equipo (\texttt{fichas\_IA\_datos\_y\_monitoreo.md}). \\
\hline
Prioridad & Media \\
\hline
\end{tabular}
\end{center}
\end{minipage}

\begin{minipage}{\linewidth}
\bigskip
\noindent\textbf{RNF-26. Explicabilidad de la clasificación de imágenes (IA-02)}
\label{tab:RNF-26}

\begin{center}
\begin{tabular}{|p{4.2cm}|p{10cm}|}
\hline
Identificador & RNF-26 \\
\hline
Nombre & Explicabilidad de la clasificación de imágenes (IA-02) \\
\hline
Característica ISO/IEC 25010:2023 & Interacción de usuario \\
\hline
Subcaracterística & Autodescriptividad \\
\hline
Descripción & El sistema deberá presentar, junto con cada propuesta de clase de RF-29, el nivel de confianza y la región de la fotografía que motivó la clasificación, de forma simultánea a la propuesta y en lenguaje comprensible para el usuario. \\
\hline
Métrica & Proporción de propuestas que incluyen clase, confianza y región señalada de forma simultánea \\
\hline
Valor objetivo & 100~\%. \\
\hline
Método de medición & Verificado sobre una muestra de 20 propuestas de clasificación. \\
\hline
Origen & RF-29; marco de explicabilidad de la Sección~\ref{sec:explicabilidad-ia} (Chazette et al.~\cite{chazette2020,chazette2021,chazette2022}); EV-01 [09:02]. \\
\hline
Prioridad & Alta \\
\hline
\end{tabular}
\end{center}
\end{minipage}

\subsection{Componentes de inteligencia artificial del sistema}
\label{sec:componentes-ia}

Conforme a la Sección 4.4 y el Paso 3 de la guía de la PE5, esta sección reúne para cada componente de inteligencia artificial los cinco bloques exigidos (descripción; datos de entrenamiento; métricas de éxito; requisitos éticos, de privacidad y clasificación de riesgo; plan de monitoreo) en una ficha única, con las trazas hacia los RF y RNF ya especificados en las secciones anteriores. El SGA incorpora dos componentes cuyo comportamiento se obtiene por aprendizaje a partir de datos y no por programación explícita: IA-01, ya presente en la línea base (RF-15, RF-26, RF-28; RNF-15 a RNF-20), y IA-02, incorporado en esta práctica (RF-29 a RF-31; RNF-21 a RNF-26).

\textbf{Justificación de IA-02 con evidencia de campo.} El segundo componente elegido es un detector de plagas y enfermedades por imagen. La evidencia disponible en este ERS respalda la elección de forma directa: el reporte de pérdidas por plagas es, explícitamente, el que el administrador señala como el más importante de analizar (EV-01 [08:17]--[08:29], origen ya citado en RF-19); el catálogo de plagas y enfermedades (RF-20) y su registro asociado (RF-16) nacen de la misma evidencia (EV-01 [01:16], [01:25]; EV-02 [01:42]); y el trabajador agrícola depende de reportar incidencias fitosanitarias al administrador (EV-02 [01:34]), hoy mediante reconocimiento visual no asistido. El detector se integra con estos artefactos ya especificados (alimenta RF-16 a través de RF-29/RF-30 y refuerza el mecanismo de alerta de RF-25 a través de RF-31) sin introducir funcionalidad no solicitada por el dominio. No se dispone en este repositorio de una segunda alternativa evidenciada con el mismo nivel de detalle (por ejemplo, un predictor de rendimiento de cosecha), por lo que la elección se sostiene en la evidencia ya trazada a RF-16, RF-19, RF-20 y RF-25 antes que en una comparación entre candidatos.

\subsubsection{Ficha del componente IA-01: alertas y recomendaciones sobre registros históricos}
\label{sec:ficha-ia01}

\begin{center}
\begin{longtable}{|p{4.2cm}|p{10cm}|}
\hline
Identificador y nombre & IA-01: Alertas y recomendaciones sobre registros históricos \\
\hline
Tarea y tipo & Recomendación y detección de anomalías (patrones) sobre los registros históricos de producción, actividades, inventario y enfermedades. \\
\hline
Entradas y salidas & Entradas: registros históricos del SGA (RF-04, RF-05, RF-06, RF-16). Salidas: recomendaciones (RF-15) y alertas tempranas (RF-26) con nivel de confianza (D-03), sujetas a confirmación humana (RF-28, R-05). \\
\hline
Datos de entrenamiento & Origen: registros propios del sistema (producción RF-05, labores RF-04, enfermedades RF-16, inventario RF-06, entregas RF-22); para la puesta en marcha, digitalización controlada de la bitácora en papel del administrador (EV-01 [02:04], [02:17]). Sin datos de terceros. Volumen: mínimo operativo de 500 registros productivos por ciclo anual; por debajo de ese volumen el componente opera solo con reglas (RF-25), sin capa predictiva. Variables: parcela, cultivo, variedad, fecha, tipo de labor, cantidad cosechada con unidad (RNF-07), enfermedad con gravedad, inventario disponible; sin datos biométricos ni de localización de personas. Etiquetado: cada alerta o recomendación queda etiquetada por la confirmación o el descarte del administrador (RF-28); sin costo adicional de anotación manual. Sesgos conocidos: subregistro de enfermedades leves en parcelas revisadas con menor frecuencia, desbalance entre parcelas pequeñas y grandes, estacionalidad climática (mitigación: RNF-20 y evaluación estratificada por mes). Base legal (LOPDP): los registros de labores contienen datos personales del trabajador agrícola (identificador, actividad, parcela); tratamiento amparado en la ejecución de la relación laboral y el interés legítimo de gestión agrícola (Art.~8 LOPDP), con finalidad declarada de planificación productiva, no de evaluación individual. \\
\hline
Métricas de éxito & Exactitud de la detección temprana $\geq$ 70~\% (RNF-16); tiempo de respuesta $\leq$ 5,0 s (RNF-17); adicionalmente, precisión de recomendaciones aceptadas $\geq$ 80~\% y tasa de falsos positivos de alerta $\leq$ 30~\%, ambas medidas sobre el registro de decisiones de RF-28. \\
\hline
Equidad & Brecha de exactitud entre cultivos de cacao y verde $\leq$ 10 puntos porcentuales (RNF-19); brecha de tasa de falsos positivos entre parcelas de alto y bajo volumen de registro $\leq$ 5 puntos porcentuales (RNF-20). \\
\hline
Explicabilidad & RNF-15: explicación causal y por evidencia, con incertidumbre explícita, simultánea a la salida, de menor densidad para el trabajador agrícola (Sección~\ref{sec:explicabilidad-ia}). \\
\hline
Riesgos y \emph{fallback} & Consecuencia del error: una verificación en campo innecesaria o una alerta omitida, nunca una acción automática (R-05, RF-28). Si el módulo falla o no está disponible, el sistema continúa operando sin recomendaciones ni alertas tempranas, sin bloquear ninguna función de registro o consulta (RNF-18). \\
\hline
Clasificación de riesgo & Riesgo mínimo conforme al Reglamento (UE) 2024/1689 (Sección~\ref{sec:riesgo-ia}). \\
\hline
Plan de monitoreo & Indicadores: tasa de aceptación de recomendaciones, proporción de alertas descartadas, cobertura de datos por parcela. Frecuencia mensual sobre el registro de decisiones (RF-28). Deriva: caída $>$10 puntos porcentuales en la exactitud respecto del mes anterior, o aceptación $<$60~\% durante dos meses, dispara reentrenamiento; si no se recupera el umbral de RNF-16, el componente vuelve al modo solo reglas (RF-25). \\
\hline
Trazas & RF-04, RF-05, RF-06, RF-15, RF-16, RF-22, RF-25, RF-26, RF-28; RNF-02, RNF-07, RNF-15 a RNF-20; decisiones D-01 a D-03; R-05. \\
\hline
\end{longtable}
\end{center}

\subsubsection{Ficha del componente IA-02: detector de plagas y enfermedades por imagen}
\label{sec:ficha-ia02}

\begin{center}
\begin{longtable}{|p{4.2cm}|p{10cm}|}
\hline
Identificador y nombre & IA-02: Detector de plagas y enfermedades por imagen \\
\hline
Tarea y tipo & Clasificación de imágenes (visión por computador): propone una clase del catálogo de plagas y enfermedades del cultivo correspondiente (RF-20), o «sin hallazgo» / «no concluyente». \\
\hline
Entradas y salidas & Entrada: fotografía tomada en campo durante el registro de una incidencia (RF-16), con cultivo y parcela. Salida: clase propuesta del catálogo (RF-20), nivel de confianza y región de la imagen que motivó la propuesta (RF-29, RNF-26). \\
\hline
Datos de entrenamiento & Origen: capturas fotográficas tomadas en campo durante RF-16, complementadas con repositorios públicos etiquetados como semilla inicial de las clases del catálogo (RF-20); las fotos de entorno ya levantadas como evidencia sirven de calibración de condiciones reales de captura. Volumen: entre 800 y 1200 imágenes etiquetadas para el arranque, mínimo 150 por clase objetivo del catálogo. Etiquetado: doble (identificación preliminar desde la semilla pública y confirmación agronómica local sobre cada imagen, desacuerdo resuelto por tercera revisión), con concordancia mínima kappa de 0,80 antes de aceptar el conjunto. Sesgos conocidos: iluminación y etapa fenológica concentradas en pocas jornadas de captura, predominio de un dispositivo móvil, síntomas avanzados sobrerrepresentados (mitigación: RNF-24, RNF-25, muestreo estratificado). Base legal (LOPDP): las imágenes de plantas no constituyen datos personales; protocolo de captura con exclusión de personas, documentos o placas del encuadre, con recorte si una persona aparece igualmente. \\
\hline
Métricas de éxito & F1 macro $\geq$ 0,85, ninguna clase por debajo de 0,75 (RNF-21); tiempo de inferencia $\leq$ 2,0 s (RNF-22); adicionalmente, degradación aceptable fuera de distribución F1 $\geq$ 0,75 en evaluación semestral previa a temporada de lluvias. \\
\hline
Equidad & Brecha de F1 entre condición de luz (plena/sombra) $\leq$ 5 puntos porcentuales (RNF-24); brecha de F1 entre gama de dispositivo $\leq$ 5 puntos porcentuales (RNF-25). \\
\hline
Explicabilidad & RNF-26: clase, confianza y región de la imagen que motivó la clasificación, simultáneas a la propuesta, comprensibles para administrador y trabajador agrícola. \\
\hline
Riesgos y \emph{fallback} & Sin conectividad o con confianza insuficiente, la propuesta se presenta como sugerencia editable y el registro continúa por selección manual del catálogo (RF-20, RNF-23); nunca bloquea el registro de la incidencia (RF-16). El registro solo se completa tras confirmación, corrección o descarte humano (RF-30, R-05). \\
\hline
Clasificación de riesgo & Riesgo mínimo conforme al Reglamento (UE) 2024/1689 (Sección~\ref{sec:riesgo-ia}). \\
\hline
Plan de monitoreo & Indicadores: concordancia entre clase propuesta y clase finalmente registrada (RF-30), proporción de capturas descartadas, distribución de condiciones de captura. Frecuencia mensual sobre los registros confirmados. Deriva: concordancia $<$80~\% o F1 $<$0,80 dispara reentrenamiento trimestral programado o inmediato; dos ciclos sin recuperar el umbral retiran la propuesta automática y el flujo vuelve a selección manual pura. \\
\hline
Trazas & RF-16, RF-19, RF-20, RF-25, RF-29, RF-30, RF-31; RNF-21 a RNF-26; decisiones D-04, D-05; R-05. \\
\hline
\end{longtable}
\end{center}

\subsubsection{Clasificación de riesgo y marco legal aplicable}
\label{sec:riesgo-ia}

Conforme al Paso 3.d de la guía, el sistema se clasifica siguiendo el enfoque basado en riesgo del Reglamento (UE) 2024/1689, que distingue prácticas prohibidas (Art.~5), sistemas de alto riesgo (Anexo~III, con obligaciones reforzadas) y obligaciones de transparencia cuando el sistema interactúa directamente con personas o genera contenido que podría confundirse con una decisión humana. Ambos componentes de IA del SGA se clasifican como de \textbf{riesgo mínimo}, por tres razones:

\begin{enumerate}
\item \textbf{No constituyen práctica prohibida.} Ninguno de los dos componentes manipula subliminalmente al usuario, asigna una puntuación social ni realiza categorización biométrica: IA-01 recomienda sobre registros agrícolas y IA-02 clasifica fotografías de plantas.
\item \textbf{No encajan en los ámbitos de alto riesgo del Anexo~III.} Las salidas de ambos componentes se refieren a cultivos y parcelas, no a personas: no deciden sobre empleo, evaluación de trabajadores, crédito, educación, servicios esenciales, identificación biométrica ni infraestructura crítica. El sistema sí gestiona información de trabajadores (RF-14, RF-01), pero esa gestión es una función determinista del SGA, no una decisión de los componentes de IA, que no evalúan el desempeño de personas ni asignan tareas.
\item \textbf{La restricción R-05 refuerza la ausencia de decisión autónoma.} Ninguna salida de IA-01 o IA-02 es accionable por sí sola: RF-28 exige confirmación o descarte de cada alerta o recomendación, y RF-30 exige confirmación, corrección o descarte de cada clasificación antes de registrarla.
\end{enumerate}

Las obligaciones que el sistema sí asume, propias de un componente de riesgo mínimo que interactúa con el usuario, son de transparencia: informar que las alertas, recomendaciones y clasificaciones son generadas por inteligencia artificial (cubierto por RNF-15 y RNF-26), mantener supervisión humana efectiva (R-05, RF-28, RF-30) y documentar técnicamente los componentes (fichas de las Secciones~\ref{sec:ficha-ia01} y~\ref{sec:ficha-ia02}).

En cuanto a la normativa ecuatoriana de protección de datos personales (Ley Orgánica de Protección de Datos Personales y su Reglamento General), el sistema en su conjunto sí trata datos personales de los trabajadores a través de RF-01 (control de acceso por rol) y RF-14 (gestión de usuarios), con cifrado de credenciales (RNF-11). Para IA-01 específicamente, el tratamiento de los registros históricos que la alimentan se ampara en la ejecución de la relación laboral y el interés legítimo de gestión agrícola (Art.~8 LOPDP), con la finalidad declarada de planificación productiva y sin uso para evaluación individual del trabajador. Para IA-02, las fotografías de plantas no constituyen datos personales bajo el protocolo de captura descrito en la ficha de la Sección~\ref{sec:ficha-ia02}, por lo que este componente no añade un tratamiento de datos personales adicional al ya existente en el sistema.

\subsection{Historias de usuario}
\label{sec:historias-usuario}

Esta sección presenta las historias de usuario en formato Connextra para cada requisito funcional de prioridad Debe tener, con su verificación según los criterios INVEST y al menos un escenario de aceptación en notación Gherkin. Los roles utilizados son Administrador y Jornalero, conforme al perfil evidenciado en las entrevistas EV-01 y EV-02. No se redactan historias de usuario para los requisitos de prioridad Debería tener, Podría tener o No tendrá, conforme al alcance definido para este bloque. El identificador de cada historia coincide con el del requisito funcional del que deriva.

\bigskip
\noindent\textbf{HU-01 \textperiodcentered\ deriva de RF-01}

\smallskip
\noindent Como usuario del sistema (administrador o jornalero), quiero iniciar sesión con mi usuario y contraseña, para acceder únicamente a las funciones correspondientes a mi rol, conforme a la restricción de seguridad R-02.

\smallskip
\noindent\textit{Verificación INVEST}
\begin{itemize}
\item \textbf{Independiente}: no depende de otra historia de usuario para implementarse.
\item \textbf{Negociable}: el mecanismo de recuperación de contraseña puede definirse en diseño.
\item \textbf{Valiosa}: sin ella ningún otro módulo puede diferenciar administrador de jornalero; origen normativo (R-02), sin evidencia de entrevista posible.
\item \textbf{Estimable}: alcance acotado a autenticación y redirección por rol.
\item \textbf{Pequeña}: cabe en una iteración.
\item \textbf{Verificable}: cumple si el sistema autentica y concede acceso según el rol en $\leq$ 2,0 s (RF-01).
\end{itemize}

\smallskip
\noindent\textit{Escenarios de aceptación (Gherkin)}

\noindent CA-01.1
\begin{itemize}
\item[] \textbf{Dado} que el usuario está registrado en el sistema
\item[] \textbf{Cuando} ingresa un usuario y una contraseña válidos
\item[] \textbf{Entonces} el sistema concede acceso a las funciones correspondientes a su rol
\end{itemize}

\bigskip
\noindent\textbf{HU-02 \textperiodcentered\ deriva de RF-02}

\smallskip
\noindent Como administrador, quiero registrar y consultar las parcelas de la finca, para contar con un registro estructurado que complemente la identificación visual que utiliza actualmente.

\smallskip
\noindent\textit{Verificación INVEST}
\begin{itemize}
\item \textbf{Independiente}: no depende de otra historia de usuario.
\item \textbf{Negociable}: los campos exactos de ubicación pueden ajustarse en diseño.
\item \textbf{Valiosa}: complementa con un registro estructurado la identificación visual que el administrador ya domina (EV-01 [04:02]).
\item \textbf{Estimable}: alcance acotado a un CRUD de parcelas.
\item \textbf{Pequeña}: cabe en una iteración.
\item \textbf{Verificable}: cumple si la consulta de una parcela responde en $\leq$ 2,0 s con 500 parcelas cargadas (RF-02).
\end{itemize}

\smallskip
\noindent\textit{Escenarios de aceptación (Gherkin)}

\noindent CA-02.1
\begin{itemize}
\item[] \textbf{Dado} que el administrador inició sesión
\item[] \textbf{Cuando} registra una parcela con nombre, ubicación, área y estado
\item[] \textbf{Entonces} la parcela queda disponible para su consulta y para asociarle cultivos
\end{itemize}

\bigskip
\noindent\textbf{HU-03 \textperiodcentered\ deriva de RF-03}

\smallskip
\noindent Como administrador, quiero registrar los cultivos de verde y cacao de cada parcela, para registrar en el sistema los cultivos correspondientes a la distribución real de la finca entre palma, verde y cacao.

\smallskip
\noindent\textit{Verificación INVEST}
\begin{itemize}
\item \textbf{Independiente}: depende únicamente de que exista una parcela (RF-02), no de otra historia de su mismo nivel.
\item \textbf{Negociable}: la lista de variedades puede ampliarse en diseño.
\item \textbf{Valiosa}: sin esto no se puede diferenciar verde de cacao en ningún reporte posterior (EV-01 [03:18]).
\item \textbf{Estimable}: alcance acotado a un CRUD de cultivos asociados a parcela.
\item \textbf{Pequeña}: cabe en una iteración.
\item \textbf{Verificable}: cumple si el cultivo queda disponible para consulta asociado a su parcela en $\leq$ 2,0 s (RF-03).
\end{itemize}

\smallskip
\noindent\textit{Escenarios de aceptación (Gherkin)}

\noindent CA-03.1
\begin{itemize}
\item[] \textbf{Dado} que existe al menos una parcela registrada
\item[] \textbf{Cuando} el administrador registra un cultivo de tipo verde o cacao para esa parcela
\item[] \textbf{Entonces} el cultivo queda asociado correctamente a la parcela
\end{itemize}

\bigskip
\noindent\textbf{HU-04 \textperiodcentered\ deriva de RF-04}

\smallskip
\noindent Como jornalero, quiero registrar las labores que realizo cada día, para dejar un registro de las actividades ejecutadas, además de la comunicación verbal que mantengo con el administrador.

\smallskip
\noindent\textit{Verificación INVEST}
\begin{itemize}
\item \textbf{Independiente}: depende de que el trabajador tenga una cuenta registrada (RF-14), no de otra historia de su mismo nivel.
\item \textbf{Negociable}: el catálogo exacto de tipos de actividad puede ajustarse en diseño.
\item \textbf{Valiosa}: añade un registro consultable a la comunicación verbal que hoy predomina (EV-02 [00:54]).
\item \textbf{Estimable}: alcance acotado a un formulario de registro.
\item \textbf{Pequeña}: cabe en una iteración.
\item \textbf{Verificable}: cumple si el 100~\% de las actividades quedan disponibles para consulta del administrador en máximo 1 minuto (RF-04).
\end{itemize}

\smallskip
\noindent\textit{Escenarios de aceptación (Gherkin)}

\noindent CA-04.1
\begin{itemize}
\item[] \textbf{Dado} que el jornalero inició sesión
\item[] \textbf{Cuando} registra una actividad indicando fecha, tipo y parcela
\item[] \textbf{Entonces} la actividad queda disponible para consulta del administrador
\end{itemize}

\bigskip
\noindent\textbf{HU-05 \textperiodcentered\ deriva de RF-05}

\smallskip
\noindent Como jornalero, quiero registrar la cantidad cosechada en la unidad que corresponde a cada cultivo, para mantener en el sistema la misma unidad de medida que utilizo actualmente según se trate de plátano o cacao.

\smallskip
\noindent\textit{Verificación INVEST}
\begin{itemize}
\item \textbf{Independiente}: depende de que exista un cultivo registrado (RF-03), no de otra historia de su mismo nivel.
\item \textbf{Negociable}: la lista de unidades puede ampliarse en diseño (RNF-14).
\item \textbf{Valiosa}: es la base de todo reporte de producción posterior (EV-02 [02:08]).
\item \textbf{Estimable}: alcance acotado a un formulario de registro de cosecha.
\item \textbf{Pequeña}: cabe en una iteración.
\item \textbf{Verificable}: cumple si el registro se incorpora al historial de producción en $\leq$ 2,0 s (RF-05).
\end{itemize}

\smallskip
\noindent\textit{Escenarios de aceptación (Gherkin)}

\noindent CA-05.1
\begin{itemize}
\item[] \textbf{Dado} que existe un cultivo previamente registrado
\item[] \textbf{Cuando} el jornalero registra una cosecha con cultivo, parcela, cantidad, unidad y fecha
\item[] \textbf{Entonces} la producción queda registrada y disponible en el historial
\end{itemize}

\bigskip
\noindent\textbf{HU-06 \textperiodcentered\ deriva de RF-06}

\smallskip
\noindent Como administrador, quiero registrar y consultar el inventario de insumos y herramientas, para dejar de operar sin un inventario real.

\smallskip
\noindent\textit{Verificación INVEST}
\begin{itemize}
\item \textbf{Independiente}: no depende de otra historia de usuario.
\item \textbf{Negociable}: los tipos de insumo pueden ampliarse en diseño.
\item \textbf{Valiosa}: resuelve la ausencia de inventario reconocida por el propio administrador (EV-01 [06:14]).
\item \textbf{Estimable}: alcance acotado a un CRUD de insumos.
\item \textbf{Pequeña}: cabe en una iteración.
\item \textbf{Verificable}: cumple si el inventario refleja la actualización en $\leq$ 2,0 s (RF-06).
\end{itemize}

\smallskip
\noindent\textit{Escenarios de aceptación (Gherkin)}

\noindent CA-06.1
\begin{itemize}
\item[] \textbf{Dado} que el administrador inició sesión
\item[] \textbf{Cuando} registra o actualiza un insumo con nombre, tipo, cantidad y unidad
\item[] \textbf{Entonces} el inventario queda actualizado con la información registrada
\end{itemize}

\bigskip
\noindent\textbf{HU-07 \textperiodcentered\ deriva de RF-07}

\smallskip
\noindent Como administrador, quiero consultar la producción registrada en un período específico, para saber si la diferencia entre un mes y otro es más alta o más baja.

\smallskip
\noindent\textit{Verificación INVEST}
\begin{itemize}
\item \textbf{Independiente}: depende de que existan registros de cosecha (RF-05), no de otra historia de su mismo nivel.
\item \textbf{Negociable}: los filtros adicionales pueden ampliarse en diseño.
\item \textbf{Valiosa}: apoya la comparación de rendimiento entre periodos (EV-01 [04:47]).
\item \textbf{Estimable}: alcance acotado a una consulta filtrada por fecha.
\item \textbf{Pequeña}: cabe en una iteración.
\item \textbf{Verificable}: cumple si la consulta responde en $\leq$ 2,0 s con 500 registros (RF-07).
\end{itemize}

\smallskip
\noindent\textit{Escenarios de aceptación (Gherkin)}

\noindent CA-07.1
\begin{itemize}
\item[] \textbf{Dado} que existen registros de producción almacenados
\item[] \textbf{Cuando} el administrador selecciona un período de consulta
\item[] \textbf{Entonces} el sistema presenta únicamente la producción de ese período
\end{itemize}

\bigskip
\noindent\textbf{HU-08 \textperiodcentered\ deriva de RF-08}

\smallskip
\noindent Como administrador, quiero planificar las actividades agrícolas con anticipación, para mantener la práctica de organizar el trabajo con una semana de adelanto.

\smallskip
\noindent\textit{Verificación INVEST}
\begin{itemize}
\item \textbf{Independiente}: depende de que existan parcelas (RF-02), no de otra historia de su mismo nivel.
\item \textbf{Negociable}: el horizonte de planificación puede ajustarse en diseño.
\item \textbf{Valiosa}: digitaliza una práctica de planificación que ya existe pero es informal (EV-01 [06:51]).
\item \textbf{Estimable}: alcance acotado a un formulario de planificación.
\item \textbf{Pequeña}: cabe en una iteración.
\item \textbf{Verificable}: cumple si la actividad planificada queda disponible para asignación en $\leq$ 2,0 s (RF-08).
\end{itemize}

\smallskip
\noindent\textit{Escenarios de aceptación (Gherkin)}

\noindent CA-08.1
\begin{itemize}
\item[] \textbf{Dado} que existen parcelas registradas
\item[] \textbf{Cuando} el administrador registra una actividad con tarea, fecha programada y parcela
\item[] \textbf{Entonces} la actividad queda disponible para su posterior asignación
\end{itemize}

\bigskip
\noindent\textbf{HU-09 \textperiodcentered\ deriva de RF-09}

\smallskip
\noindent Como administrador, quiero asignar tareas específicas a cada trabajador, para dejar un registro formal de la planificación y dirección del trabajo que realizo directamente con cada uno.

\smallskip
\noindent\textit{Verificación INVEST}
\begin{itemize}
\item \textbf{Independiente}: depende de RF-14 y RF-08, no de otra historia de su mismo nivel.
\item \textbf{Negociable}: el nivel de detalle de la tarea puede ajustarse en diseño.
\item \textbf{Valiosa}: deja registro formal de la coordinación que hoy realiza directamente el administrador (EV-02 [01:01]).
\item \textbf{Estimable}: alcance acotado a una asignación trabajador-tarea-parcela-fecha.
\item \textbf{Pequeña}: cabe en una iteración.
\item \textbf{Verificable}: cumple si el 100~\% de las tareas asignadas quedan disponibles para el trabajador en $\leq$ 2,0 s (RF-09).
\end{itemize}

\smallskip
\noindent\textit{Escenarios de aceptación (Gherkin)}

\noindent CA-09.1
\begin{itemize}
\item[] \textbf{Dado} que existen trabajadores y actividades registradas
\item[] \textbf{Cuando} el administrador asigna una actividad a un trabajador con parcela y fecha
\item[] \textbf{Entonces} la tarea queda asignada y disponible para consulta del trabajador
\end{itemize}

\bigskip
\noindent\textbf{HU-10 \textperiodcentered\ deriva de RF-10}

\smallskip
\noindent Como administrador, quiero consultar el estado de las tareas asignadas, para contar con un registro consultable del seguimiento que realizo actualmente de forma directa con cada trabajador.

\smallskip
\noindent\textit{Verificación INVEST}
\begin{itemize}
\item \textbf{Independiente}: depende de RF-09, no de otra historia de su mismo nivel.
\item \textbf{Negociable}: los estados intermedios pueden ampliarse en diseño.
\item \textbf{Valiosa}: deja un registro consultable del seguimiento que hoy realiza el administrador de forma directa (EV-01 [07:11]).
\item \textbf{Estimable}: alcance acotado a una consulta de estado.
\item \textbf{Pequeña}: cabe en una iteración.
\item \textbf{Verificable}: cumple si el estado se refleja correctamente en $\leq$ 2,0 s desde el cambio (RF-10).
\end{itemize}

\smallskip
\noindent\textit{Escenarios de aceptación (Gherkin)}

\noindent CA-10.1
\begin{itemize}
\item[] \textbf{Dado} que existen tareas previamente asignadas
\item[] \textbf{Cuando} el administrador consulta el módulo de seguimiento
\item[] \textbf{Entonces} el sistema muestra el estado actualizado de cada tarea asignada
\end{itemize}

\bigskip
\noindent\textbf{HU-11 \textperiodcentered\ deriva de RF-11}

\smallskip
\noindent Como administrador, quiero generar reportes de producción por cultivo, parcela o período, para analizar el rendimiento de la finca sin calcularlo manualmente.

\smallskip
\noindent\textit{Verificación INVEST}
\begin{itemize}
\item \textbf{Independiente}: depende de que existan registros de producción, no de otra historia de su mismo nivel.
\item \textbf{Negociable}: el formato de exportación puede ajustarse en diseño.
\item \textbf{Valiosa}: responde al tipo de reporte que el administrador identificó como necesario (EV-01 [08:10]).
\item \textbf{Estimable}: alcance acotado a un generador de reportes con filtros.
\item \textbf{Pequeña}: cabe en una iteración.
\item \textbf{Verificable}: cumple si el reporte se genera en $\leq$ 3,0 s con 500 registros (RF-11).
\end{itemize}

\smallskip
\noindent\textit{Escenarios de aceptación (Gherkin)}

\noindent CA-11.1
\begin{itemize}
\item[] \textbf{Dado} que existen registros de producción almacenados
\item[] \textbf{Cuando} el administrador selecciona los filtros de cultivo, parcela o período
\item[] \textbf{Entonces} el sistema genera el reporte correspondiente disponible para consulta o descarga
\end{itemize}

\bigskip
\noindent\textbf{HU-14 \textperiodcentered\ deriva de RF-14}

\smallskip
\noindent Como administrador, quiero registrar y administrar las cuentas de los usuarios del sistema con su rol correspondiente, para administrar el acceso del equipo de siete personas conforme al rol de cada una.

\smallskip
\noindent\textit{Verificación INVEST}
\begin{itemize}
\item \textbf{Independiente}: no depende de otra historia de usuario.
\item \textbf{Negociable}: los campos adicionales del perfil pueden ampliarse en diseño.
\item \textbf{Valiosa}: sostiene el control de acceso diferenciado del que depende todo el sistema (EV-01 [02:35]; R-02).
\item \textbf{Estimable}: alcance acotado a un CRUD de usuarios y roles.
\item \textbf{Pequeña}: cabe en una iteración.
\item \textbf{Verificable}: cumple si el cambio de rol o estado se refleja en los permisos en $\leq$ 2,0 s (RF-14).
\end{itemize}

\smallskip
\noindent\textit{Escenarios de aceptación (Gherkin)}

\noindent CA-14.1
\begin{itemize}
\item[] \textbf{Dado} que el administrador inició sesión
\item[] \textbf{Cuando} registra un usuario con su rol correspondiente
\item[] \textbf{Entonces} la cuenta queda disponible con los permisos de ese rol
\end{itemize}

\bigskip
\noindent\textbf{HU-19 \textperiodcentered\ deriva de RF-19}

\smallskip
\noindent Como administrador, quiero un reporte de las pérdidas de producción por plagas y enfermedades, para tener la información que considero la más importante de analizar.

\smallskip
\noindent\textit{Verificación INVEST}
\begin{itemize}
\item \textbf{Independiente}: depende de que existan registros de enfermedad (RF-16), no de otra historia de su mismo nivel.
\item \textbf{Negociable}: el desglose exacto (por cultivo, por parcela) puede ajustarse en diseño.
\item \textbf{Valiosa}: responde al reporte que el propio administrador señaló como el más importante (EV-01 [08:17]).
\item \textbf{Estimable}: alcance acotado a un reporte con filtros.
\item \textbf{Pequeña}: cabe en una iteración.
\item \textbf{Verificable}: cumple si el reporte se genera en $\leq$ 3,0 s con 500 registros (RF-19).
\end{itemize}

\smallskip
\noindent\textit{Escenarios de aceptación (Gherkin)}

\noindent CA-19.1
\begin{itemize}
\item[] \textbf{Dado} que existen registros de enfermedad con pérdida estimada
\item[] \textbf{Cuando} el administrador solicita el reporte de pérdidas por plagas
\item[] \textbf{Entonces} el sistema presenta la pérdida estimada por plaga, cultivo y parcela
\end{itemize}

\bigskip
\noindent\textbf{HU-24 \textperiodcentered\ deriva de RF-24}

\smallskip
\noindent Como administrador, quiero consultar el historial de labores y cosechas filtrado por trabajador y fecha, para no volver a perder el control de los datos como ocurrió cuando no se llevaba bitácora de las fundas semanales.

\smallskip
\noindent\textit{Verificación INVEST}
\begin{itemize}
\item \textbf{Independiente}: depende de que existan registros de RF-04 y RF-05, no de otra historia de su mismo nivel.
\item \textbf{Negociable}: los filtros adicionales pueden ampliarse en diseño.
\item \textbf{Valiosa}: previene la desviación de datos ya ocurrida una vez (EV-01 [02:04], [02:17]).
\item \textbf{Estimable}: alcance acotado a una consulta filtrada.
\item \textbf{Pequeña}: cabe en una iteración.
\item \textbf{Verificable}: cumple si la consulta con 500 registros responde en $\leq$ 2,0 s (RF-24).
\end{itemize}

\smallskip
\noindent\textit{Escenarios de aceptación (Gherkin)}

\noindent CA-24.1
\begin{itemize}
\item[] \textbf{Dado} que existen labores y cosechas registradas
\item[] \textbf{Cuando} el administrador filtra el historial por trabajador y rango de fechas
\item[] \textbf{Entonces} el sistema presenta únicamente los registros que cumplen ese filtro
\end{itemize}

\bigskip
\noindent\textbf{HU-25 \textperiodcentered\ deriva de RF-25}

\smallskip
\noindent Como administrador, quiero recibir alertas automáticas cuando algo requiera atención inmediata, para tomar decisiones sin tener que revisar todo el sistema manualmente.

\smallskip
\noindent\textit{Verificación INVEST}
\begin{itemize}
\item \textbf{Independiente}: depende de RF-06 y RF-16 como origen de la alerta, no de otra historia de su mismo nivel.
\item \textbf{Negociable}: el canal de notificación puede definirse en diseño.
\item \textbf{Valiosa}: es la forma de apoyo a la decisión que el administrador priorizó explícitamente (EV-01 [10:33]).
\item \textbf{Estimable}: alcance acotado a un motor de alertas sobre condiciones ya definidas.
\item \textbf{Pequeña}: cabe en una iteración.
\item \textbf{Verificable}: cumple si el 100~\% de las enfermedades de gravedad alta generan alerta en $\leq$ 5,0 s (RF-25).
\end{itemize}

\smallskip
\noindent\textit{Escenarios de aceptación (Gherkin)}

\noindent CA-25.1
\begin{itemize}
\item[] \textbf{Dado} que se registra una enfermedad con nivel de gravedad alto
\item[] \textbf{Cuando} el sistema procesa el registro
\item[] \textbf{Entonces} genera una alerta visible al administrador en menos de 5 segundos
\end{itemize}

\bigskip
\noindent\textbf{HU-26 \textperiodcentered\ deriva de RF-26}

\smallskip
\noindent Como administrador, quiero que el sistema me avise automáticamente cuando una parcela muestre señales de enfermedad antes de que yo lo note visualmente, para intervenir a tiempo.

\smallskip
\noindent\textit{Verificación INVEST}
\begin{itemize}
\item \textbf{Independiente}: depende de que existan registros históricos suficientes, no de otra historia de su mismo nivel.
\item \textbf{Negociable}: el algoritmo específico de detección puede definirse en diseño técnico.
\item \textbf{Valiosa}: es la funcionalidad de IA con la que el administrador se mostró explícitamente de acuerdo (EV-01 [09:02]).
\item \textbf{Estimable}: alcance acotado a la generación de una alerta temprana con nivel de confianza.
\item \textbf{Pequeña}: cabe en una iteración, asumiendo el modelo de IA como componente ya disponible.
\item \textbf{Verificable}: cumple si el módulo alcanza una exactitud mínima del 70~\% sobre el conjunto de validación (RF-26).
\end{itemize}

\smallskip
\noindent\textit{Escenarios de aceptación (Gherkin)}

\noindent CA-26.1
\begin{itemize}
\item[] \textbf{Dado} que existen registros históricos suficientes de una parcela
\item[] \textbf{Cuando} el módulo de IA detecta un patrón asociado a una posible enfermedad
\item[] \textbf{Entonces} el sistema genera una alerta temprana con el nivel de confianza del modelo
\end{itemize}

% Bibliografia
\bibliographystyle{ieeetr}
\bibliography{referencias}

\end{document}
